% ---------------------------------------------------------------------
% explog.tex -- OptEvolve experiment log and working notes.
%
% A standalone companion to submission.tex, not part of it.  The paper
% cites sections here through the `xr' package, and this document cites
% the paper the same way, so BOTH must be compiled twice for the
% cross-document numbers to settle:
%
%     tectonic submission.tex && tectonic explog.tex
%     tectonic submission.tex && tectonic explog.tex
%
% Style rules: never use the em-dash; no italic text; -ize/-ization
% spelling (customization, factorization, optimization), except -wise
% suffixes, revise, exercise, noise/denoising, and file paths, which keep
% their own spelling.
% ---------------------------------------------------------------------
\documentclass[11pt]{article}

\usepackage[margin=2.4cm]{geometry}
\usepackage[utf8]{inputenc}
\usepackage[T1]{fontenc}
\usepackage{times}
\usepackage{amsmath,amsfonts,amssymb}
\usepackage{booktabs}
\usepackage{graphicx}
\usepackage{xcolor}
\usepackage{nicefrac}
\usepackage{microtype}
\usepackage{placeins}
\usepackage{natbib}
\usepackage{xr-hyper}
\usepackage[colorlinks=true,linkcolor=blue!45!black,
            citecolor=blue!45!black,urlcolor=blue!45!black]{hyperref}

% Labels of the paper are reachable here as PAPER-<label>. Keep
% submission.pdf beside this PDF for the links to open the paper.
\externaldocument[PAPER-]{submission}[submission.pdf]

\title{\textbf{OptEvolve: Experiment Log and Working Notes}\\[2mm]
\large Companion record to the paper\\
\large OptEvolve: Hardware-Aware Discovery of Certified Optimization Algorithms}
\date{}

\begin{document}
\maketitle
\label{app:log}

This appendix is the running record of the project. It lists what was
measured, on which data and hardware, and what was tried and abandoned.
Negative results are kept on purpose: several of them decide what the paper
may and may not claim. Every number is traceable to an artifact in the code
repository; the path is given in each subsection.

\tableofcontents
\clearpage

\section{Three axes of customization}
\label{app:log:axes}

A solver is varied along three axes, and each experiment below moves one of
them. Naming them separately is what makes the negative results readable: a
recipe that wins on one family and loses on another has not failed, it has
located an axis-two boundary.

\begin{enumerate}
\item The solving algorithm. The base operator, the wrappers composed onto it
 (relaxation, anchoring, fixed and adaptive restarts, safeguarded
 acceleration), the metric, and the certificate cadence. This axis acts
 mainly on $N$, the number of iterations needed to certify.
\item The input problem family. Real data at several scales, described by the
 structure of the instance and by its conditioning. This axis does not make a
 recipe faster; it decides which recipe is right, and which baseline is even
 able to run.
\item The hardware. Execution structure, storage precision, and kernel and
 iteration fusion. This axis acts on $t$, the cost of one iteration.
\end{enumerate}

The product $T = N \times t$ is what a user experiences, and the certificate
cadence is the only gene that appears in both factors, which is why it is
treated separately throughout.

\section{What was run, and what was found}
\label{app:log:scoreboard}

Every row is measured against solvers written by their authors, on downloaded
real data except where marked generated, with acceptance being our own
float64 certificate in the original coordinates. Detail and artifact paths are
in the subsections that follow.

\begin{center}
\small
\begin{tabular}{p{2.5cm}p{2.6cm}p{2.2cm}p{6.1cm}}
\toprule
Family & Data and scale & Baseline & Outcome \\
\midrule
TV denoising & six images, $n=128$ to $496$ & our own XLA loop & Fusion wins on the hardware axis: $1.36\times$ to $2.63\times$ end to end on the RTX 4090, $1.15\times$ to $1.73\times$ on the Blackwell. Admission is bitwise identical at $\rho=1$. \\
\addlinespace
Elastic-net logistic & epsilon (LIBSVM), $10^5\times2000$, ten penalty settings & cuML, same GPU & $2.28\times$ ($95\%$ interval $[1.73, 3.01]$) on ten of ten for a repeated solve, and $0.06\times$ for a first solve on an empty cache. The margin is monotone in the penalty weight and crosses over near $0.05\lambda_{\max}$. \\
\addlinespace
Convex QP, Maros-Meszaros & 137 published instances & OSQP, SCS, Clarabel & No win. We certify 55 to 63 instances against OSQP's 123 and Clarabel's 132, and are $4.5\times$ to $5.8\times$ slower where both solve. Every instance factorizes cheaply. \\
\addlinespace
Convex QP, generated lasso & $n$ from $25{,}000$ to $750{,}000$ variables & OSQP, SCS, Clarabel & $1.98\times$ over Clarabel at $25{,}000$ and $6.04\times$ at $75{,}000$. At $250{,}000$ we finish in $102.7$ s while Clarabel is killed by the operating system at $28.9$ GB and OSQP does not certify in $5400$ s. Our own cost per iteration also degrades with size. \\
\addlinespace
Convex QP, generated SVM & $n$ from $3300$ to $110{,}000$ & OSQP, SCS, Clarabel & The ADMM solvers are beaten from about $11{,}000$ variables, and Clarabel never is: its lead grows with size, because the coupling constraints give an interior-point method a system that factorizes well. Measured before the timing correction described above and not repeated. \\
\addlinespace
Optimal transport, DOTmark & $32\times32$ and $64\times64$ & POT exact network simplex & Lost by $5\times$ to $32\times$ and abandoned. The exact simplex is very strong at these sizes, and the entropic scheme never certified on the original problem. \\
\bottomrule
\end{tabular}
\end{center}

Read along axis two, the pattern is that the wins are family dependent rather
than merely size dependent. A large instance is not sufficient: the generated
SVM family is large and is never won against Clarabel, while the generated
lasso family of the same size is, because the two give an interior-point
method very different systems to factorize. The one clean statement at the
largest sizes is not about speed at all, but about which solvers still run.

No new algorithm is claimed. The ingredient that carries the logistic result,
recomputing the Lipschitz constant on the current working set, is published:
Wang, Yu, Liang, Wu and Yu, An Iterative Reduction FISTA Algorithm for
Large-Scale LASSO, SIAM Journal on Scientific Computing 44(4), 2022, use it on
the same epsilon data at the same penalty weight and attribute their speedup
to the same smaller constant. What is ours is the measurement, the attribution
between ingredients, and the derivation of the recipe from measured symptoms.

\section{Completed-field evaluation, September 16}
\label{app:field}

The tables above compare our solvers with generic solvers only. This subsection
reports what happens when the field is completed: the structure-specific solver
for each family is added where one exists, every competitor is swept over a
ladder of its own tolerance settings, and every returned point is judged by our
acceptance check on the original problem. Artifacts are in
\texttt{results/lasso\_specialists\_20260915}, \texttt{results/field\_audit\_20260916},
\texttt{results/completed\_field\_20260916}, \texttt{results/lasso\_path\_cv\_20260916},
\texttt{results/preconditioning\_20260916}, \texttt{results/dotmark\_ot\_20260915}
and \texttt{results/synthesis\_transfer\_20260915}.

\paragraph{Specialists on the families of the earlier tables.}
For lasso the specialist runs on the native data $(A, b, \lambda)$, recovered
from the QP encoding and checked to reproduce the QP objective exactly, and is
scored at the fastest setting whose original objective is within $10^{-6}$
relative of the best objective any arm reached.

\begin{center}
\small
\begin{tabular}{p{2.1cm}p{2.3cm}rrp{2.6cm}r}
\toprule
family & instance & ours & Clarabel & specialist & time \\
\midrule
lasso & $n = 10{,}000$ & 2.395\,s & 1.960\,s & scikit-learn & 0.0018\,s \\
lasso & $n = 50{,}000$ & 4.168\,s & 181.7\,s & scikit-learn & 0.0199\,s \\
lasso & $n = 100{,}000$ & 20.52\,s & not run & scikit-learn & 0.0479\,s \\
huber & $n_{\mathrm{var}} = 140{,}000$ & 2.068\,s & 132.38\,s & L-BFGS-B, unlifted & 0.267\,s \\
MEB dual & $n_{\mathrm{var}} = 165{,}000$ & 3.35--3.62\,s & 135--202\,s & Frank--Wolfe, away steps & 0.437\,s \\
group lasso & $n = 10{,}000$ & 3.645\,s$^{a}$ & 14.92\,s & celer & 0.674\,s \\
SVM & $n_{\mathrm{var}} = 110{,}000$ & 10.35\,s & 39.8--43.7\,s & LIBLINEAR & 43.0\,s \\
transport & DOTmark $32\times32$ & 328.9\,s$^{b}$ & 255.6\,s & network simplex & 4.797\,s \\
\bottomrule
\end{tabular}
\end{center}

Lasso is at tolerance $10^{-4}$ for our arm and Clarabel on the RTX 4090, with
our time the minimum over paired repetitions; the $n = 100{,}000$ figure is a
re-measurement (20.517\,s and 22.279\,s) that replaces an unarchived 17.604\,s.
Huber, the MEB dual and SVM are at $10^{-6}$; group lasso and transport at
$10^{-3}$. $^{a}$First visit of the routed system, including exploration.
$^{b}$Our best formulation, the standard form; the relocated form cuts
iterations from $91{,}450$ to $29{,}800$ but raises the per-iteration cost from
$3.43\times10^{-4}$\,s to $2.62\times10^{-2}$\,s and takes $891.9$\,s. At the
transport size, HiGHS needs $118.1$\,s at a setting our check accepts. The
lasso solutions have $91$, $1{,}136$ and $2{,}807$ nonzeros, so coordinate
descent works on a tiny fraction of the columns. On control at $n_x = 1{,}200$,
OSQP is $1.59\times$ faster than Clarabel at $10^{-8}$ with objectives agreeing
to nine digits ($2.57\times$ and $3.03\times$ at $n_x = 600$ and $300$), so the
control margin of Table~\ref{PAPER-tab:general} is measured against a generic solver
that is not the strongest one on that family. On SVM the specialist does not
beat our method. Every specialist named here is a documented, published method
\citep{Pedregosa2011Scikit,Massias2018Celer,Byrd1995Limited,Yildirim2008Two,LacosteJulien2015Global,Fan2008LIBLINEAR,Bonneel2011Displacement,Flamary2021POT};
none was discovered by us.

\paragraph{The held-out field, completed.}
The held-out routing comparison of Section~\ref{PAPER-sec:exp:specialists} was rerun
with ECOS offered second-order cone programs (it had been installed but never
offered one, and it is the only engine that reaches the target on
\texttt{CBLIB:qssp60}), and with recognition-gated specialists for network
flow, group lasso and the minimum enclosing ball. The minimum enclosing ball
specialist was written by us from the published Frank--Wolfe method, not taken
from a package. Tolerance $10^{-3}$, 20\,s budget, RTX 4090.

\begin{center}
\scriptsize
\begin{tabular}{lrrrll}
\toprule
instance & routed repeat (s) & first visit (s) & oracle (s) & oracle engine & verdict \\
\midrule
H:knapsack\_ridge:12000:0.003 & 0.2057 & 1.0070 & 6.5223 & clarabel & routed 31.71$\times$ \\
H:knapsack\_ridge:48000:0.0005 & 0.9902 & 2.8940 & -- & -- & routed only \\
H:svm\_dual:6000 & 0.1309 & 0.4860 & 0.0478 & piqp & oracle 2.74$\times$ \\
H:svm\_dual:20000 & 0.2890 & 0.8920 & 1.3485 & clarabel & routed 4.67$\times$ \\
H:portfolio:20000 & 0.0890 & 0.5710 & 0.0442 & piqp & oracle 2.01$\times$ \\
H:portfolio:100000 & 0.2387 & 1.1970 & 0.3404 & piqp & routed 1.43$\times$ \\
H:budget\_band:20000 & 0.1004 & 0.4080 & 0.0831 & piqp & oracle 1.21$\times$ \\
H:budget\_band:80000 & 0.1925 & 1.0340 & 0.5542 & piqp & routed 2.88$\times$ \\
MM:DUALC5 & 0.0014 & 0.4920 & 0.0011 & piqp & oracle 1.26$\times$ \\
MM:DUALC8 & 0.0020 & 0.3900 & 0.0016 & piqp & oracle 1.21$\times$ \\
MM:HS118 & 0.0007 & 1.9410 & 0.0007 & piqp & routed 1.09$\times$ \\
MM:QAFIRO & 0.0007 & 0.4390 & 0.0008 & piqp & routed 1.07$\times$ \\
HLP:transport:60 & 0.0004 & 0.4770 & 0.0007 & flow:pot & routed 1.51$\times$ \\
HLP:mcf\_joint:12 & 0.0139 & 1.7990 & 0.0081 & piqp & oracle 1.71$\times$ \\
HLP:gub\_packing:2560 & 0.0834 & 2.2350 & 0.0322 & piqp & oracle 2.60$\times$ \\
HLP:production:1536 & 0.1281 & 1.7180 & 0.0645 & piqp & oracle 1.99$\times$ \\
HLP:netflow\_cap:48 & 0.0947 & 0.2650 & 0.0262 & piqp & oracle 3.61$\times$ \\
NETLIB:AFIRO & 0.0007 & 0.4220 & 0.0007 & piqp & oracle 1.02$\times$ \\
NETLIB:SC105 & 0.0011 & 2.0880 & 0.0010 & piqp & oracle 1.15$\times$ \\
NETLIB:ADLITTLE & 0.0011 & 1.8090 & 0.0009 & piqp & oracle 1.16$\times$ \\
HSOCP:group\_lasso:10000 & 0.2784 & 3.6450 & 10.6844 & glasso:skglm & routed 38.38$\times$ \\
HSOCP:portfolio\_risk:50000 & 1.9850 & 4.3100 & 4.8391 & scs & routed 2.44$\times$ \\
HSOCP:facility\_location:10000 & 0.0558 & 5.4870 & 0.0922 & clarabel & routed 1.65$\times$ \\
HSOCP:robust\_lp:5000 & 0.4483 & 5.9950 & -- & -- & routed only \\
HSOCP:min\_enclosing\_ball:10000 & 0.4288 & 3.9210 & 0.6356 & meb:coreset & routed 1.48$\times$ \\
CBLIB:nb & 0.4124 & 0.9740 & 0.4210 & scs & routed 1.02$\times$ \\
CBLIB:qssp60 & -- & -- & 1.3762 & ecos & oracle only \\
CBLIB:nql30 & 0.0524 & 0.9200 & 0.0579 & clarabel & routed 1.11$\times$ \\
\bottomrule
\end{tabular}
\end{center}

Both sides solve $25$ instances; routing is faster on $13$ and the oracle on
$12$; routing alone solves $2$ and the oracle alone $1$. The instance
\texttt{H:knapsack\_ridge:48000:0.0005} was solved in $2$ of $5$ reruns at this
budget. \texttt{CBLIB:qssp60} is lost to exploration order rather than to the
field: ECOS is among the router's candidates but is seventh of seven, and the
budget expires after six. The group lasso first visit became slower
($0.71$\,s to $3.65$\,s) because the recognition-gated specialist is explored
before our cone formulation, which is faster on that instance; we report this
rather than reorder the candidates on held-out data. The oracle for group lasso
in this table is skglm at $10.68$\,s; a later sweep with celer added gives
$0.674$\,s, so the group lasso row is not a win.

\paragraph{The chosen routes, timed head to head.}
The routing runs above report a first visit that includes exploration and a
repeat visit predicted from calibration. We timed the route the system chose on
six of the instances directly. Our arm runs in a fresh process with structure
detection, formulation build, compilation and solve inside one clock (cold), and
a second solve follows in the same process (warm); five repetitions, RTX 4090,
both arms pinned to the same sixteen cores, every run on a quiet machine.
Competitors run in fresh processes over the tolerance ladder with a 600\,s limit.
Ratios are competitor time over ours, paired geometric means with $95\%$
intervals; references for knapsack\_ridge are certified by a duality bound
(relative gap at most $8\times10^{-11}$). Artifacts are in
\texttt{results/framework\_best\_20260916}.

\begin{center}
\scriptsize
\begin{tabular}{lrrrr}
\toprule
instance & tolerance & ours cold / warm (s) & best engine (s) & speedup, cold \\
\midrule
knapsack\_ridge, $n = 12{,}000$ & $10^{-3}$ & 1.04 / 0.48 & Clarabel 4.22 & $4.08\times$ $[3.92, 4.25]$ \\
 & $10^{-4}$ & 1.04 / 0.51 & Clarabel 4.69 & $4.49\times$ $[4.45, 4.53]$ \\
 & $10^{-6}$ & 1.13 / 0.60 & Clarabel 5.66 & $5.00\times$ $[4.93, 5.07]$ \\
knapsack\_ridge, $n = 48{,}000$ & $10^{-3}$ & 6.42 / 2.43 & Clarabel 224.7 & $34.98\times$ $[34.61, 35.36]$ \\
 & $10^{-4}$ & 8.10 / 2.81 & Clarabel 204.3 & $25.22\times$ $[24.95, 25.49]$ \\
 & $10^{-6}$ & 11.91 / 3.75 & Clarabel 245.0 & $20.57\times$ $[20.39, 20.74]$ \\
robust LP, $n = 5{,}000$ & $10^{-3}$ & 1.18 / 0.84 & Clarabel 14.89 & $12.64\times$ $[12.26, 13.02]$ \\
 & $10^{-4}$ & 1.77 / 1.38 & Clarabel 15.71 & $8.86\times$ $[8.33, 9.43]$ \\
 & $10^{-6}$ & 2.97 / 2.63 & Clarabel 16.75 & $5.63\times$ $[5.18, 6.12]$ \\
SVM dual, $n = 20{,}000$ & $10^{-3}$ & 1.14 / 0.58 & Clarabel 0.95 & $0.84\times$ $[0.82, 0.86]$ \\
 & $10^{-4}$ & 1.20 / 0.63 & Clarabel 1.24 & $1.03\times$ $[1.01, 1.05]$ \\
 & $10^{-6}$ & 1.31 / 0.71 & Clarabel 1.44 & $1.11\times$ $[1.07, 1.14]$ \\
budget band, $n = 80{,}000$ & $10^{-3}$ & 1.31 / 0.76 & PIQP 0.31 & $0.24\times$ $[0.23, 0.25]$ \\
 & $10^{-4}$ & 1.33 / 0.78 & PIQP 0.32 & $0.24\times$ $[0.23, 0.25]$ \\
 & $10^{-6}$ & 1.38 / 0.81 & PIQP 0.34 & $0.25\times$ $[0.24, 0.25]$ \\
group lasso, $n = 10{,}000$ & $10^{-3}$ & 1.00 / 0.63 & celer 0.57 & $0.57\times$ $[0.56, 0.57]$ \\
 & $10^{-4}$ & 1.02 / 0.64 & celer 0.57 & $0.56\times$ $[0.55, 0.56]$ \\
 & $10^{-6}$ & 1.04 / 0.66 & celer 0.57 & $0.55\times$ $[0.54, 0.56]$ \\
\bottomrule
\end{tabular}
\end{center}

On group lasso, skglm takes $7.7$\,s with its numba compilation in the clock and
$1.64$\,s without it. On the SVM dual, LIBSVM's SMO solves the identical dual and
passes our check in $19$ to $21$\,s; LIBLINEAR was not run because it does not
enforce the dual equality constraint, so its answer has no exact lift. The
knapsack\_ridge $n = 12{,}000$ competitor is slower in the routing run
($6.52$\,s) than here ($4.22$\,s) because six engines ran at once there.

\paragraph{Where the speedup comes from.}
Holding the relocated formulation fixed and swapping only the projection, on
knapsack\_ridge at $n = 12{,}000$ and tolerance $10^{-4}$ (cold / warm, seconds):

\begin{center}
\small
\begin{tabular}{lrrr}
\toprule
projection & cold & warm & per-iteration share \\
\midrule
bisection, 60 steps & 1.65 & 1.07 & $97\%$ \\
breakpoint search (textbook) & 4.35 & 0.72 & $93\%$ \\
Newton, 12 steps (hand-written) & 1.26 & 0.68 & $92\%$ \\
synthesized, round 1 & 1.04 & 0.53 & $79\%$ \\
synthesized, round 2 & 1.05 & 0.52 & $71\%$ \\
synthesized in the later agent session (best) & 1.05 & 0.51 & $66\%$ \\
\bottomrule
\end{tabular}
\end{center}

Against the best engine's $4.69$\,s, relocation with the hand-written Newton
projection is already $3.72\times$ faster, and the synthesized operator adds
$1.21\times$ cold and $1.32\times$ warm on top. The textbook breakpoint search is
slow cold because its sort takes about $3$\,s to compile. Although the
projection is $66\%$ to $97\%$ of each iteration, the projection's total time is
only $1.4\%$ to $7.3\%$ of a cold first solve on the four relocated instances
measured:
detection ($0.13$ to $1.2$\,s), build and compilation dominate. On the budget
band the operator does not matter, because the solve takes only $60$ to $260$
iterations.

\paragraph{Scale.}
On knapsack\_ridge with $\rho = 0.003$ at tolerance $10^{-4}$, with the best
synthesized operator:

\begin{center}
\small
\begin{tabular}{rrrrrr}
\toprule
$n$ & ours, cold (s) & breakpoint, cold (s) & best engine (s) & speedup, cold & warm \\
\midrule
12{,}000 & 1.05 & 4.30 & 4.8 & $4.58\times$ $[4.48, 4.68]$ & $9.32\times$ \\
24{,}000 & 2.83 & 6.65 & 32.2 & $11.37\times$ $[11.18, 11.58]$ & $27.73\times$ \\
48{,}000 & 5.09 & 9.99 & 219.2 & $43.09\times$ $[41.82, 44.40]$ & $99.48\times$ \\
100{,}000 & 25.6 & 28.6 & no answer & over $23\times$ & over $25\times$ \\
\bottomrule
\end{tabular}
\end{center}

The margin grows with $n$ because the engines' factorizations grow faster than
our iteration cost; the synthesized operator's advantage over the breakpoint
search shrinks with $n$ ($4.1\times$ to $1.12\times$ cold). At $n = 100{,}000$
Clarabel exceeded a 24\,GB memory limit and PIQP returned nothing in $600$\,s, so
the entry is a lower bound. Our own time grows $5\times$ from $48{,}000$ to
$100{,}000$ because the router's adaptive chunk sizing recompiles the iteration
$12$ to $20$ times; with fixed $50$-iteration chunks the $n = 100{,}000$ solve
takes $9.6$\,s cold and the held-out $n = 48{,}000$ solve at $10^{-4}$ takes
$3.3$\,s, which is not what the router currently ships.

\paragraph{A compilation cache that is not cold.}
The package enables a persistent on-disk compilation cache. A new process
therefore reloads kernels compiled by earlier processes, and a first visit timed
in a fresh process in the routing runs above is not a cold solve. The three tables
above were measured with the cache disabled.

\paragraph{Tolerance settings are not a common currency.}
With two settings per competitor, $176$ rejections over $125$
(instance, solver, target) combinations were soft: the solver returned a point
whose violation was still falling toward our target when the ladder ended.
Deepening to four settings on both sides, ours included, turned $14$ of the
$54$ eligible combinations into acceptances ($26\%$), raised the accepted
records from $427$ to $455$, and left the oracle's solve count unchanged at
$26$ of $28$. The HiGHS rejections are not of this kind: all of them hit the
time limit, and their objective gap is identical at every setting. The same
error affected a lasso path comparison of our own. Run at the setting whose
numeral matched the target, celer and scikit-learn appeared to certify only
$75$ and $52$ of $100$ regularization values at $10^{-8}$, because their
stopping tests normalize the duality gap by the objective at zero rather than
at the returned point. Swept to $10^{-10}$ through $10^{-14}$ they certify all
$100$ and are faster than our batched GPU method:

\begin{center}
\small
\begin{tabular}{lrrr}
\toprule
100-value lasso path, all values at $10^{-8}$ & scikit-learn & celer & ours (H100, float64) \\
\midrule
$n = 10{,}000$ & 1.03\,s & not recorded & 2.00\,s \\
$n = 50{,}000$ & 4.97\,s & 5.91\,s & 7.57\,s \\
$n = 100{,}000$ & 11.46\,s & 18.70\,s & 23.09\,s \\
\bottomrule
\end{tabular}
\end{center}

\paragraph{Batching, cross-validation and many right-hand sides.}
A batched proximal-gradient step applies one sparse product to all path values
at once, but that product does $\mathrm{nnz}(A) \cdot L$ work for $L$ values, and
only the index traffic and the kernel launch are shared, so batching buys at
most about $1.4\times$, exhausted by $L = 100$. The specialists, by contrast, get
cheaper per value as the grid grows denser, because warm starts get closer:
the marginal cost falls from $12.5$\,ms to $1.0$\,ms as $L$ goes from $25$ to
$3{,}200$, and the gap to our method widens $17\times$ with width. $K$-fold
cross-validation, with a row mask per fold on one shared matrix, all cells
certified on both sides at $10^{-8}$:

\begin{center}
\small
\begin{tabular}{rrrrr}
\toprule
$n$ & $K$ & ours & best specialist & ratio \\
\midrule
10{,}000 & 5 & 13.12\,s & 1.95\,s & $0.15\times$ \\
10{,}000 & 10 & 27.44\,s & 2.30\,s & $0.08\times$ \\
50{,}000 & 5 & 76.88\,s & 5.30\,s & $0.07\times$ \\
50{,}000 & 10 & 104.92\,s & 7.16\,s & $0.07\times$ \\
100{,}000 & 5 & 175.35\,s & 15.79\,s & $0.09\times$ \\
100{,}000 & 10 & 310.62\,s & 18.62\,s & $0.06\times$ \\
\bottomrule
\end{tabular}
\end{center}

This is not an implementation defect. Our sparse kernel sustains $564$ to
$726$\,GB/s of gathered traffic, $17\%$ to $22\%$ of the H100's memory
bandwidth; granting both a kernel at the bandwidth ceiling and perfect removal
of the $69\%$ to $86\%$ of column iterations spent on already-certified values,
the best case is $1.51\times$ at $n = 10{,}000$, $K = 5$ and a loss at every
larger case ($0.65\times$ at $n = 100{,}000$, $K = 10$). On $200$ unrelated
right-hand sides, which removes the specialists' warm starts entirely, we lose
to a single CPU core ($0.24\times$ at $n = 10{,}000$), and on the elastic-net
path we lose at $0.37\times$ to $0.56\times$.

\paragraph{Where the dense method wins.}
The crossover is solution density, not path width. Extending the path to
smaller $\lambda$ grows the active set until coordinate descent has no small
working set left:

\begin{center}
\small
\begin{tabular}{lllrr}
\toprule
$n$ & $\lambda_{\min}/\lambda_{\max}$ & tail support & speedup at $10^{-4}$ & at $10^{-6}$ \\
\midrule
10{,}000 & $10^{-2}$ & $17\%$ of columns & $0.52\times$ & $0.82\times$ \\
10{,}000 & $10^{-3}$ & $41\%$ of columns & $1.52\times$ & $0.66\times$ \\
10{,}000 & $10^{-4}$ & $49\%$ of columns, $98\%$ of rows & $15.8\times$ & $5.07\times$ \\
50{,}000 & $10^{-4}$ & $49\%$ of columns, $98\%$ of rows & $7.34\times$ & ours $100$/$100$ in $177$\,s \\
\bottomrule
\end{tabular}
\end{center}

At $n = 50{,}000$ and $10^{-6}$ the best completed specialist run is celer at
$72$ of $100$ values in $641$\,s; one scikit-learn setting was stopped by its job
time limit, so that cell is not fully measured. At $10^{-8}$ neither side
completes at either size.

\paragraph{Preconditioning, measured with seeds.}
Earlier figures comparing relocation, row scaling and an optimized diagonal
preconditioner were single-seed iteration counts. We repeated the comparison on
$48$ instances (eight families, two sizes, three seeds, three paired
repetitions, tolerance $10^{-6}$, H100), with wall time as the primary number
and the construction of every preconditioner inside the clock:

\begin{center}
\small
\begin{tabular}{lrr}
\toprule
pair, over instances where both finish & wall-time ratio, $95\%$ interval & iterations \\
\midrule
optimized diagonal / Ruiz & $6.28\times$ $[4.79, 8.43]$ slower & $0.47\times$ \\
relocation / optimized diagonal & $0.185\times$ $[0.087, 0.398]$ faster & $0.08\times$ \\
relocation / Ruiz & $4.10\times$ $[2.16, 7.48]$ slower & $0.09\times$ \\
optimized diagonal, construction excluded / Ruiz & $0.709\times$ $[0.532, 0.936]$ faster & $0.40\times$ \\
Ruiz / none & $0.944\times$ $[0.892, 0.997]$ faster & $0.81\times$ \\
\bottomrule
\end{tabular}
\end{center}

The optimized diagonal reduces iterations on $13$ instances and turns that
into a wall-time win on none; its construction is a median $87\%$ of its wall
time. The relocation-versus-Ruiz ratio covers only the $27$ of $48$ pairs where Ruiz
finished inside the $200{,}000$-iteration cap. Where it did not (every seed of
the MEB dual at both sizes, and portfolio and budget\_band at the larger size)
relocation takes about $0.4$\,s; where it did, relocation still wins on portfolio,
budget\_band and the SVM dual, and loses on the three LP families and on
knapsack\_ridge, whose projection costs more per iteration than the iterations
it saves. With
its construction removed the optimized diagonal is worth $1.41\times$ over Ruiz,
and it pays for itself only after a median of $17$ solves of the same matrix
(quartiles $11$ and $46$). Its construction cost grows with exponent $2.07$,
which projects to $226$ hours and a $74.5$\,GiB dense matrix at
$n = 100{,}000$. Iteration counts vary by up to $72.7\times$ across seeds of one
family and size.

\paragraph{A published learned solver under the same protocol.}
Applying the shifted geometric mean of \citet{Ichnowski2021Accelerating}, with
shift $10$, to the per-problem tables of their arXiv appendix reproduces their
solved counts on Maros--M\'esz\'aros exactly ($126$, $125$, $127$). The reported $1.829\times$ lies
between the iteration ratios over the $119$ problems OSQP and the best RLQP
policy both solve ($1.853\times$) and the $110$ all three arms solve
($1.798\times$), while the wall-clock ratio over the $119$ is $1.552\times$. Over all $12$ QPLIB instances with timeouts at the stated
$300$\,s the wall-clock ratio is $0.972\times$, and over the $8$ both solve it is
$1.88\times$. On Netlib the both-solved ratio is $0.573\times$, and the reported
$1.30\times$ is reproduced over the $97$ problems we could parse only with a fill
of about $1000$\,s for timeouts. No trained policy weights are published; we built the solver from
source (an OSQP 0.6.2 fork) but could not run the learned policy.

\section{Convex QP against production solvers, September 13}
\label{app:log:qp2026}

The earlier QP rows were measured before three defects in the harness were
found and fixed, and they are superseded by what follows. The defects were:
baselines imported into the same process as our solver, so on a machine whose
solver environment has no baselines every baseline arm recorded an import
error and our time printed against three failures; a duplicated Ruiz
equilibration, run once by the caller to read the step sizes and again inside
the preparation path; and a dual mapping that overwrote rather than summed the
multiplier on rows with two finite bounds, which made Clarabel and SCS appear
to fail on the control family when their points were feasible to $10^{-15}$.

Baselines now run in their own interpreters, so a baseline never shares an
environment with our solver and OSQP stays on the CPU. Every repetition pays
full setup on both sides, the persistent compilation cache is disabled and the
JAX caches are cleared before each of our repetitions, so both arms are timed
in the first-solve regime. Intervals are geometric means of paired log ratios.

\begin{center}
\small
\begin{tabular}{p{1.9cm}p{1.9cm}p{2.9cm}p{2.9cm}p{3.2cm}}
\toprule
Family & Size ($n$) & Across seeds, $10^{-4}$ & Accuracy matched, $10^{-6}$ & Certificate, ours vs theirs \\
\midrule
huber & 140{,}000 & $50.71\times$ $[43.51, 59.10]$, 3 seeds & $63.89\times$ $[63.67, 64.12]$ & $6.4\mathrm{e}{-}08$ vs $1.5\mathrm{e}{-}07$ \\
\addlinespace
lasso & 100{,}000 & $13.07\times$ $[11.75, 14.55]$ at $50{,}000$, 5 seeds & $55.60\times$ $[54.30, 56.93]$ & $5.1\mathrm{e}{-}08$ vs $4.2\mathrm{e}{-}06$ \\
\addlinespace
svm & 110{,}000 & $6.20\times$ $[5.44, 7.05]$, 3 seeds & $4.17\times$ $[3.71, 4.70]$ & $9.4\mathrm{e}{-}08$ vs $6.3\mathrm{e}{-}06$ \\
\addlinespace
control & 16{,}200 & $1.63\times$ $[1.59, 1.66]$, 3 seeds & $1.896\times$ $[1.673, 2.149]$ & $5.9\mathrm{e}{-}07$ vs $3.5\mathrm{e}{-}06$ \\
\addlinespace
portfolio & all & $0.97\times$ $[0.55, 1.73]$ at rank 2500 & loses at every regime & does not certify at $100{,}316$ \\
\bottomrule
\end{tabular}
\end{center}

All ratios are against Clarabel. OSQP and SCS were not run at these sizes, and
Clarabel is not the strongest generic solver on every family: on control, OSQP
is faster than Clarabel (Appendix~\ref{app:field}), and structure-specific
solvers, which are faster still on huber and lasso, were not in this field. Four of the five are won, and in each of those our certificate is
tighter than the baseline's, so none of the four is bought by returning a
weaker point. Tightening the target from $10^{-4}$ to $10^{-6}$ makes three of
the four margins larger, because our cost is dominated by preparation that does
not depend on the target while an interior-point factorization's accuracy
requirement does. The exception is svm, whose own cost exponent in $n$ is
superlinear at $1.187$ against huber's $0.668$ and lasso's $0.726$, and which
therefore pays for the tighter target in iterations.

\paragraph{Where the time goes.} At lasso with $250{,}000$ variables one run took
$17.604$ s, of which rescaling is $12.437$ s, preparation $0.107$ s,
compilation $3.973$ s and the iteration loop $1.087$ s. That run's artifact was
not archived; a re-measurement on the RTX 4090 on September 16 took $20.517$ s
and $22.279$ s, and that measurement is the one used in Appendix~\ref{app:field}. The loop that the
factorization $T = N \times t$ describes is six per cent of the runtime.
Across every family and size measured, preparation is between sixty and ninety
seven per cent of it. Two changes follow directly and account for most of the
improvement over the earlier rows: removing the duplicated equilibration, and
capping the power iteration at one hundred steps rather than five hundred,
which moves the operator norm by $9\mathrm{e}{-}10$ relative against a bound
that is inflated by $1.02$ in any case.

\paragraph{Placement is a measured decision, not a default.} The same svm
instance loses at $0.41\times$ on the H100 and wins at $1.22\times$ on the RTX
4090 with nothing else changed. The H100 has the weaker single thread, and the
work that dominates our runtime is single threaded numpy and scipy. Measured
device records give $79.3$ and $1.357$ TFLOPS on the 4090 against $362.2$ and
$52.5$ on the H100 in float32 and float64, so a decision taken from peak
throughput would choose the H100 at every size and be wrong here. The lasso
win reproduces on both, at $13.38\times$ $[13.01, 13.75]$ and $7.26\times$
$[6.10, 8.65]$.

\paragraph{Two negative results.} First, no operator chain found by the search
beats plain Condat-Vu at scale. On lasso at $50{,}000$ the base recipe needs
$14{,}080$ evaluations; Halpern with an adaptive restart certifies but is
$1.76\times$ slower, Anderson with memory five certifies but is $2.64\times$
slower, and Halpern without a restart does not certify within $200{,}000$
evaluations. Anderson's evaluation saving, which is $4.8\times$ at $1{,}500$
variables, is $1.4$ per cent at $50{,}000$. Second, no structural screen
predicts which families are winnable. Fill ratio ranks control last and
portfolio fourth, inverting the two that decide the question; factorization
work per variable, fitted post hoc on five families, then fails two of four
within a single family at fixed size. Both are reported because both were
stated in advance and refuted by measurement, and because the second failure
shows that a screen modelling only the baseline's factorization cost cannot
work when the crossing is driven from both sides.

\section{Conventions}
\label{app:log:conventions}

\begin{itemize}
\item Hardware. An NVIDIA RTX 4090 (sm\_89, local, shared with an LLM
 inference server) and an NVIDIA RTX PRO 6000 Blackwell Max-Q (sm\_120,
 97.9 GB, a shared four-GPU server of which one GPU is used for timing).
 CPU baselines run on an Intel Core i7-12700K, single threaded, one process
 pinned to one performance core.
\item Certificates. Always evaluated in float64 and in the original problem
 coordinates, whatever precision or rescaling the solver uses internally.
 LP and OT: the PDLP-style relative gap (primal residual, dual residual,
 objective gap). QP: OSQP's primal and dual residuals plus the duality gap.
\item Timing. Our solver's execution time excludes JIT compilation (reported
 separately). Baseline times include their own setup, which favours our
 solver on very small instances; this is stated wherever it matters.
\item Data. Every problem family is benchmarked on downloaded real data with
 a source URL and a sha256 manifest. Synthetic generators appear only in
 unit tests and early probes and are marked as synthetic below.
\item Instance manifests. Every iteration count names its data seed and draw
 size, because both change the sampled instances.
\end{itemize}

\section{Data provenance}
\label{app:log:data}

\begin{center}
\small
\begin{tabular}{p{2.3cm}p{4.6cm}p{5.6cm}}
\toprule
Family & Data & Source \\
\midrule
TV denoising & Kermany 2017 retinal OCT B-scans (real clinical speckle) &
 Hugging Face \\
LP & 297 instances: MIPLIB 2017 (236), fctp transportation (32), Mittelmann
 network (7); 47 above $10^6$ nonzeros & MIPLIB and Mittelmann benchmark
 pages; manifest \\
Min-cost flow & road networks (DC, DE), a vision instance (246K nodes, 1.43M
 arcs), DIMACS netgen & \url{http://lime.cs.elte.hu/~kpeter/data/mcf}
 (no reader yet) \\
Elastic net & rcv1, real-sim, news20, E2006 & LIBSVM data sets \\
Convex QP & Maros-Meszaros, 138 instances, $n$ from 2 to 93263 &
 \url{https://github.com/qpsolvers/maros_meszaros_qpbenchmark} \\
Optimal transport & DOTmark 1.0: 10 classes, resolutions 32 to 512, 500
 images & \url{https://www.stochastik.math.uni-goettingen.de/DOTmark_1.0.zip} \\
\bottomrule
\end{tabular}
\end{center}

Synthetic generators used in probes: grid and random-regular min-cost flow,
point-cloud optimal transport, and multi-commodity flow built on the grid
generator.

\section{The factorization $T = N \times t$ on TV denoising}
\label{app:log:factorisation}

Time to a certified solution factors into the iteration count $N$ and the
per-iteration cost $t$. Measured on TV denoising of $256\times256$ OCT
patches, six images, RTX 4090, tolerance $10^{-4}$:

\begin{itemize}
\item $N$ is invariant to hardware: 584, 597 and 683 at cadence $K=16, 64,
 256$, identical across two GPU architectures, both precisions, both
 execution structures, and every kernel.
\item Scheme identity is inert at coarse cadence: PDHG and Condat-Vu give 683
 and 683 at $K=256$. Step sizes are worth $4.31\times$ on $N$ and decide
 feasibility at $10^{-6}$ (the balanced setting certifies 0 of 6).
\item The precision gene was confounded with execution structure: the
 float64 path ran a device loop with a traced inner trip count, the float32
 path ran host-chunked phases with a static one. Separated, float32 buys
 nothing inside the device loop and $1.58\times$ once the loop is chunked.
\end{itemize}

\begin{center}
\small
\begin{tabular}{llccc}
\toprule
Execution & Precision & $K=16$ & $K=64$ & $K=256$ \\
\midrule
device loop & f64 & 22.54 & 20.39 & 19.78 \\
device loop & f32 & 22.99 & 20.41 & 19.78 \\
host phases & f64 & 27.18 & 16.19 & 13.29 \\
host phases & f32 & 28.00 & 12.23 & 8.43 \\
\bottomrule
\end{tabular}

Per-iteration cost $t$ in microseconds.
\end{center}

The loop penalty is one line of code: a traced inner trip count costs
$2.02\times$ per iteration at $n=256$ and $1.18\times$ at $n=496$ against a
static one (nesting and \texttt{while\_loop} are free). On the CPU the same
choice costs $1.24\times$ at $n=128$ and inverts to $0.84\times$ at $n=256$.

\section{Temporal fusion kernel}
\label{app:log:fusion}

A CUDA kernel runs $T$ PDHG or Condat-Vu iterations per launch on a
$(B+2T)\times(B+2T)$ tile with halo in shared memory. Numerical admission
against $T$ sequential reference launches: bitwise identical at $\rho=1$, at
most $4.649\times10^{-15}$ at $\rho=1.9$, on both architectures. Total wall
time over six images to a certified $10^{-4}$ gap, $K=256$, RTX 4090:

\begin{center}
\small
\begin{tabular}{rccccc}
\toprule
$n$ & baseline (s) & static loop & best fusion (tile) & total & worst tile \\
\midrule
128 & 0.06647 & $1.673\times$ & $1.571\times$ ($B=8$, $T=4$) & $2.629\times$ & $1.320\times$ \\
256 & 0.07733 & $1.347\times$ & $1.386\times$ ($B=24$, $T=4$) & $1.867\times$ & $1.298\times$ \\
384 & 0.13486 & $1.291\times$ & $1.266\times$ ($B=24$, $T=4$) & $1.634\times$ & $0.926\times$ \\
496 & 0.22564 & $1.093\times$ & $1.240\times$ ($B=32$, $T=4$) & $1.356\times$ & $0.743\times$ \\
\bottomrule
\end{tabular}
\end{center}

On the catalog's own instance draw at $n=256$, $K=256$, the gain is
$1.857\times$. This remains the only result in the project that beats software
people actually use (XLA's compiled loop).

\paragraph{The tile is a per-device decision, measured on two architectures.}
The comparison was repeated on September 17 on the RTX 4090 and an H100 80\,GB,
sweeping 21 tiles ($B \in [8, 56]$, $T \in \{2, 4, 8\}$) with the XLA static
loop and a single-step custom call on the ballot at every cell, over six OCT
images and four montage problems, to a certified $10^{-4}$ gap at $K=256$.
Both kernels build on the H100 with the stock toolchain and no source change,
and the numerical admission is identical on the two devices: bitwise identical
to $T$ sequential unfused launches at $\rho = 1$ (every cell, both devices) and
$4.635\times10^{-15}$ at $\rho = 1.9$, over 168 cases on the 4090 and 252 on
the H100 with no errors. Iteration counts are identical across arms at every
size.

\begin{center}
\small
\begin{tabular}{rllrlrr}
\toprule
 & \multicolumn{2}{c}{RTX 4090} & & \multicolumn{2}{c}{H100} & cross-device \\
\cmidrule(lr){2-3}\cmidrule(lr){5-6}
$n$ & tile & gain & & tile & gain & penalty (H100 pick on 4090) \\
\midrule
128  & $B{=}8$, $T{=}4$  & $1.660\times$ & & $B{=}16$, $T{=}8$ & $2.739\times$ & $1.304\times$ \\
256  & $B{=}24$, $T{=}4$ & $1.380\times$ & & $B{=}16$, $T{=}8$ & $2.525\times$ & $1.628\times$ \\
384  & $B{=}24$, $T{=}4$ & $1.332\times$ & & $B{=}24$, $T{=}8$ & $2.087\times$ & $1.359\times$ \\
496  & $B{=}32$, $T{=}4$ & $1.257\times$ & & $B{=}32$, $T{=}8$ & $1.947\times$ & $1.214\times$ \\
992  & $B{=}40$, $T{=}4$ & $1.146\times$ & & $B{=}40$, $T{=}4$ & $2.013\times$ & $1.000\times$ \\
1488 & $B{=}24$, $T{=}4$ & $1.373\times$ & & $B{=}40$, $T{=}4$ & $2.428\times$ & $1.009\times$ \\
\bottomrule
\end{tabular}
\end{center}

The two devices choose differently at five of the six sizes, and carrying a
choice across costs $1.06\times$ to $1.63\times$; the penalty is asymmetric,
since the 4090 choice run on the H100 costs at most $1.188\times$. The 4090
never selects $T = 8$ and the H100 selects it at every size up to $496$. Two
mechanisms separate them and both are structural rather than incidental: a
$T = 8$ tile recomputes its halo $2.25$ to $4.00$ times over, which the H100
can pay for in arithmetic and the 4090's rate-limited consumer float64 cannot;
and seven of the 21 tiles exceed the 4090's $101{,}376$-byte opt-in shared
memory per block and cannot launch there at all, against the H100's
$232{,}448$. Rebuilding with native SASS for each architecture instead of the
shipped PTX selects the same tile at every size on both devices, so the
dependence is not an artifact of just-in-time compilation.

Choosing wrongly is worse than not fusing: on the 4090, seven to eight of the
14 launchable tiles are slower than the static loop at every size, the worst at
$0.506\times$. A selector that queries the CUDA driver for the device name,
compute capability, streaming multiprocessor count and shared memory per block,
and then looks the tile up, returns $T = 4$ tiles on the 4090 and $T = 8$ tiles
on the H100 from one file with no recompile. This is table lookup conditioned
on detected hardware, not kernel generation. Its honest accuracy is the
leave-one-device-out figure, not the in-table one: applying the H100 table on
the 4090 costs a geometric mean of $1.252\times$ (worst $1.628\times$) and the
reverse $1.099\times$ (worst $1.188\times$). Interpolating across sizes on one
device fails, at zero of six exact on the 4090.

\paragraph{The decline-to-fuse branch is unexercised.}
Not fusing was a candidate at all twelve device-size cells and won none of
them, on either device. The earlier Blackwell cell at $n = 496$, where the best
tile was $0.925\times$ and the right decision was not to fuse, is the only
evidence for that branch, and it is now doubtful: that sweep's grid contained
no $T = 8$ tile, and $T = 8$ is what wins on the H100 at every size up to
$496$. The Blackwell part is unavailable to retest, so we report the branch as
machinery that measurement has not yet justified.

\section{Other problem families}
\label{app:log:families}

\paragraph{Min-cost network flow (synthetic grid and regular generators).}

The execution-structure gain replicates on a sparse LP: mean $1.386\times$
(range 1.277 to 1.492, 24 runs), with $N$ identical across loop structures
in 24 of 24 runs and identical between PDHG and Condat-Vu. The step sizes
that were merely slow on TV are inadmissible here: the construction gate
rejects them before execution ($t s \|L\|^2 = 1.02267$), 24 of 48
configurations. In the structure sweep (2106 runs, 1755 admitted, 351
rejected) admission depends only on the fill $\theta = t s \|L\|^2$: every
$\theta\le0.99$ is admitted and $\theta=1.02$ never is. The admissible
region is family dependent: $\|A\|=\sqrt{m+n}$ for transport, a region
shrinking as $\Theta(1/K)$ in the number of commodities for multi-commodity
flow, and a norm that saturates at 4 on grids.

\paragraph{The primal-weight rule.}
 A step-ratio rule
$r = C/\omega_0^p$ with $\omega_0 = \|c\|/\|b\|$ was fitted from data. It ties
PDLP's primal-weight heuristic ($p=2$): $1.0704\times$ against $1.0704\times$
of the per-instance oracle on 88 ultrawide netflow instances, $1.1822\times$
against $1.1829\times$ on 80 fine-grid instances, $1.1984\times$ against
$1.2020\times$ on 108 wide instances. The fitted exponent (about 1.8) moved
from 2.00 to 1.80 as grid censoring was removed. Transferred to transport and
multi-commodity flow, the rule reaches $1.125\times$ of the oracle on 162
held-out instances against $1.825\times$ for the best fixed ratio. We
reproduce PDLP's rule; we do not beat it.

\paragraph{Elastic net (LIBSVM).}

Coordinate descent with working sets (celer, adelie, glmnet) dominates. A
full-gradient method loses $13.6\times$ (rcv1) and $18.8\times$ (news20) core
for core. On the RTX 4090 against one CPU core it ties rcv1 (0.130 s against
0.137 s), wins news20 by $1.45\times$, and loses about $550\times$ on E2006,
which must be centered (uncentered, one near-constant feature solves it at
every regularization value tried). Verdict: a poor target, because the
in-class state of the art is not first-order.

\section{Convex QP against author code (Maros-Meszaros)}
\label{app:log:qp}

TBD.

\section{A knowledge base of classical schemes}
\label{app:log:kb}

Methods adopted in LP practice
since 2021 (restarts, Halpern reflection, averaging, primal weights) trace
back to fixed-point papers from 1950 to 1980 through the monotone-inclusion
literature. The knowledge base records, for each classical scheme, its
original paper, update, side condition, guarantee, genome location, gate
rule, and whether it has already been used in general primal-dual splitting
and in LP or QP practice (with evidence or the queries that found nothing).

\begin{center}
\small
\begin{tabular}{p{2.6cm}p{2.2cm}p{7.6cm}}
\toprule
Scheme & Origin & Status and outcome \\
\midrule
Krasnoselskii-Mann, Halpern, restarts & 1953 to 1967 & implemented (genome wrapper genes) \\
Anderson acceleration & 1965 & already used for PDHG on LP (arXiv 2508.08062); not frontier \\
Passty averaging & 1979 & already used (restarted averaged PDHG, arXiv 2206.03041); not frontier \\
Haugazeau & 1968 & implemented as an outer wrapper, with the projection evaluated in the PDHG metric; admission-tested. Rejected by measurement on QP: solves 8 and 9 instances against 55 and 63 without it, needs $20.5\times$ and $34.3\times$ more iterations where both solve, and solves nothing new \\
Bregman (entropic) & 1967 & implemented as an entropic PDHG. The gate admits it only where the constraints imply a simplex (transport), with a size-independent step condition $t s \|A\|_{1\to2}^2<1$, $\|A\|_{1\to2}=\sqrt2$. Needs a dual recovery step; never certified on DOTmark \\
Solodov-Svaiter & 1999 & frontier, not yet tried \\
Ishikawa & 1974 & not found in primal-dual practice; not tried \\
\bottomrule
\end{tabular}
\end{center}

Features harvested from author code and recorded for implementation as
combinable genes: an adaptive step ratio (OSQP's adaptive $\rho$, PDLP's
primal weight), a factorization-based prox for the quadratic (which would
make PDHG admissible on QP), and polishing (an active-set reduced solve
accepted only if the certificate improves).

\section{Optimal transport on DOTmark}
\label{app:log:ot}

 Baseline: POT's exact
network simplex (\texttt{ot.emd}) on one CPU core, 0.04 to 0.85 s per pair at
$32\times32$ (about $10^6$ arcs) and 0.8 to 44 s at $64\times64$ (about
$1.7\times10^7$ arcs).

\begin{itemize}
\item Execution structure. Writing the transport operator as row and column
 sums of the plan instead of scatter-adds over arcs (same arithmetic)
 cuts the RTX 4090's per-iteration time from 243 to 75 microseconds at
 $32\times32$ and from 6.32 to 1.07 ms at $64\times64$. One CPU core needs
 15 ms per iteration at $32\times32$ even with the faster form, so the CPU
 arm was not pursued.
\item Step ratio. A fixed ratio grid over $[0.03, 30]$, and a wider one, certified
 nothing in $5\times10^4$ iterations: the plan stayed at zero. With the ratio
 set from $\omega_0 = \|c\|/\|b\| \approx 8.8\times10^6$, plain PDHG certifies
 $10^{-4}$ in 5184 iterations. Reflected Halpern restarts did not help; the
 entropic scheme never certified.
\item Certificate weakness. The LP certificate normalizes the primal residual
 by $1+\|b\|$; for probability measures $\|b\|\approx0.01$, so a $10^{-4}$
 pass is almost an absolute bound and admitted a plan with $4.7\%$
 objective error. Transport results are therefore judged by relative error
 against the exact cost.
\item Head to head. At $32\times32$ and $10^{-6}$ (objective error
 $5\times10^{-7}$): 3.0 s on the GPU against 0.21 s for POT, $14\times$
 slower. At $64\times64$ (ClassicImages): 6.2 s against 19.4 s, but with
 $4.7\%$ objective error, which is not a solution; 20.4 s against 19.4 s at
 $7\times10^{-4}$ objective error, a tie. No win at matched meaningful
 accuracy.
\end{itemize}

\paragraph{Creation loop on transport (corrected composer).}
On CauchyDensity at resolution 32, pair 1001-1002 (tolerance $10^{-6}$, primal
weight $\omega_0/30$, one Blackwell GPU), the gates admit 19 of 48 compositions
of PDHG. Plain PDHG certifies in 69,248 evaluations and 4.27 s. The best
composition is reflected Halpern with a fixed restart every 512 steps: 61,440
evaluations and 3.90 s, $1.13\times$ fewer evaluations and $1.10\times$ faster.
With a restart every 64 steps it is roughly neutral; residual-triggered restarts
cost more evaluations or fail, and unanchored Halpern and Anderson acceleration
reach the cap of $10^5$ evaluations. A first run of this loop reported the
residual-triggered recipes as $12\times$ faster than plain PDHG; that was an
execution artifact of the composer, not a
property of the recipes.

\section{Certified adaptive step ratio}
\label{app:log:adaptive}

PDHG and Condat-Vu are proximal steps in the metric
$M = [[I/t, -L^\top], [-L, I/s]]$, so changing the primal weight
$\omega = \sqrt{s/t}$ changes the metric. With $\theta = \sqrt{ts}\,\|L\| < 1$,
$M \succeq (1-\theta)D$ for its block diagonal $D$, so a move from $(t,s)$ to
$(t',s')$ satisfies $M' \preceq (1+\eta)M$ with
$\eta = \max(0, t/t'-1, s/s'-1)/(1-\theta)$ (checked numerically over 200
random moves). The controller applies PDLP's primal-weight update every 512
iterations, places every $(t,s)$ on the gate boundary, and stops adapting
once $\sum\eta$ reaches a declared budget, which keeps the variable metric
convergence theorem of Combettes and V\~u in force. The comparison arm runs
through the same runner with the controller off.

\paragraph{QP (Maros-Meszaros, 137 instances, CPU, tolerance $10^{-4}$).}
With Ruiz rescaling the adaptive ratio passes the certificate on 73
instances against 65 with the fixed ratio (15 gained, 7 lost); where both
pass it needs $0.53\times$ the iterations (geometric mean; fewer on 35 of 58).
Without rescaling: 62 against 55 (9 gained, 2 lost), $0.50\times$ the
iterations. The budget froze 8 runs; the median run made 30 updates. Under the
accuracy audit (Section~\ref{app:log:qp}) the gain survives: with Ruiz
rescaling 59 against 50 answers are not grossly wrong (33 against 25 strict),
and without rescaling 53 (25 strict). That makes it our best QP arm, ahead of
Condat-Vu with Anderson (54), but OSQP has 89 such answers and Clarabel 116,
and where both pass the certificate it remains $7.8\times$ slower than OSQP. A
real gain inside the method, not a win.

\paragraph{Optimal transport (DOTmark $32\times32$, 31 pairs, RTX 4090, tolerance $10^{-6}$).}
Started at PDLP's initial weight $\omega_0 = \|c\|/\|b\|$, neither the fixed nor
the adaptive ratio certifies any pair within $10^5$ iterations, although
adaptation cuts the median objective error from $2.7\times10^{-2}$ to
$1.1\times10^{-3}$; it moves the ratio at every epoch and never settles. The
fixed ratio $\omega_0/30$ calibrated on a held-out pair certifies 29 of 31, a
median $19\times$ slower than POT (best $4.7\times$). At $64\times64$ nothing
certifies within $10^5$ iterations (105 s against 39 to 47 s for POT). On
transport, adaptation does not replace calibration.
Adapting for a bounded number of updates and then freezing (epochs of 2048
iterations) does not rescue it: at most 10 updates from $\omega_0$ certifies 5
of 31 pairs, at most 30 updates certifies 12, and 10 updates started from the
calibrated ratio certifies 6, against 29 for the calibrated fixed ratio. The
runs that do certify are faster (a median 11.7 to $14\times$ POT against
$21.6\times$), so the rule helps when it lands, but it does not land reliably.

\section{Elastic-net logistic regression}
\label{app:log:logreg}

Proposed by N.~Devanathan: elastic-net logistic regression on epsilon (dense,
$4\times10^5\times2000$), rcv1 (sparse), Tabula Sapiens (1.1M cells) and
ogbg-molpcba (437,929 molecules, 128 assays, multitask), against cuML's GPU
quasi-Newton (OWL-QN), skglm, and glmnet-style adelie. Two items to add in
order: fixed diagonal preconditioning followed by adaptation, and
certificate-driven structure reduction by safe screening. The possibility
the user singled out is that screening shrinks the problem mid-solve until
the best device flips (GPU first, CPU afterwards), which no existing solver
does automatically.

Built so far. Atoms for the logistic loss ($L_f = \|X\|_2^2/(4n)$) and the
elastic-net penalty (scaled soft threshold). The certificate is the relative
duality gap at the dual point $u = \nabla F(Xw)$; for $\alpha<1$ the penalty's
conjugate is smooth, for $\alpha=1$ the point is rescaled into dual
feasibility. FISTA with O'Donoghue--Cand\`es gradient restart is a new scheme
with gate $0 < a \le 1/L$ and an independent reference in the equivalence
battery. Gap Safe screening (Ndiaye, Fercoq, Gramfort and Salmon, 2017)
uses the radius $\sqrt{\mathrm{gap}/(2n)}$, since the dual is $4n$-strongly
concave (the binary-entropy term has curvature at least 4).

Defect caught. For $\alpha=1$ the bare strict screening test discarded two
active features at an exactly converged point: active features sit on
$|x_j^\top u^\star| = \lambda$, and the rescaling into dual feasibility can leave
them one rounding step below $\lambda$. A relative margin of $10^{-9}$ fixes it;
the safety test (no screened feature nonzero at a $10^{-12}$-accurate optimum)
passes for $\alpha \in \{0.5, 0.9, 1\}$.

Baseline check. On synthetic data every baseline certifies the same
objective under our gap: cuML with $C = 1/(n\lambda)$ reaches $2\times10^{-15}$,
confirming the mapping; skglm, scikit-learn saga and adelie also pass.
Baselines run in isolated environments on both machines. The training file
of epsilon is being resumed after a truncated first download.

\paragraph{First head-to-head (epsilon test split, real data).}
$n=10^5$, $p=2000$, $\alpha=0.5$, $\lambda=0.1\lambda_{\max}$, target relative duality
gap $10^{-6}$ on our certificate; GPU runs on the RTX 4090, CPU baselines on
one core. 

\begin{center}
\small
\begin{tabular}{llrrl}
\toprule
Solver & device & seconds & iterations & gap \\
\midrule
cuML (OWL-QN, float32) & 4090 & 0.183 & 35 & $2.2\times10^{-7}$ \\
ours, FISTA (float64) & 4090 & 0.668 & 192 & $6.3\times10^{-9}$ \\
ours, FISTA (float32 iterates, float64 certificate) & 4090 & 0.722 & 192 & $5.9\times10^{-9}$ \\
ours, forward-backward & 4090 & 4.88 & 1408 & $7.0\times10^{-7}$ \\
adelie (30-point path) & 1 core & 9.05 & & $2.2\times10^{-8}$ \\
scikit-learn saga & 1 core & 21.8 & & $8.5\times10^{-11}$ \\
skglm (tolerance tightened to $10^{-6}$) & 1 core & 59.2 & & $3.5\times10^{-8}$ \\
\bottomrule
\end{tabular}
\end{center}

FISTA on the GPU is $13.5\times$ faster than the fastest single-core baseline,
which is not a like-for-like hardware comparison. On the same GPU, cuML is
$3.7\times$ faster than our FISTA. Float32 iterates buy nothing (0.722 s against
0.668 s). The iteration is expected to be memory bound, since the design is
1.6 GB in float64 and is read twice per step (for $Xw$ and for $X^\top(\cdot)$);
a fused single-pass gradient kernel that reads each row block once for both
products is the natural next step. The first cuML attempt ran out of GPU
memory because our process still held the data next to the language-model
server; it was rerun alone.

\paragraph{Where the step's time goes.}
 In float64 the
step is memory bound at about 90\% of peak bandwidth: $Xw$ takes 1.74 ms (918
GB/s), $X^\top v$ 1.84 ms (871 GB/s), and the gradient 3.52 ms, two full reads
of $X$, which is essentially the whole measured FISTA iteration (3.48 ms). In
float32, $Xw$ is bandwidth bound and twice as fast (0.89 ms), but $X^\top v$
runs at only 298 GB/s (2.70 ms), so the float32 gradient is no faster than
the float64 one. Storing $X^\top$ explicitly, or writing the product as $vX$,
changes nothing (294 to 296 GB/s): the loss is the library kernel XLA selects
for this shape (2000 outputs, each a sum of $10^5$ terms), not the memory
layout. Mixed precision therefore pays only with a kernel that reaches full
bandwidth on the transposed product; a fused single pass in float32 would
cost about 0.9 ms per gradient, four times less than float64, and bfloat16
storage would halve that again. Low-precision data changes the problem, so
reaching a $10^{-6}$ gap on the true problem needs a precision schedule that
returns to float64 gradients as the gap closes; the certificate stays float64
throughout.

\begin{center}
\small
\begin{tabular}{lrrl}
\toprule
Form & ms & GB/s & relative error \\
\midrule
$X^\top v$, as shipped & 2.589 & 309 & $2.9\times10^{-4}$ \\
contracting with a float64 accumulator & 1.417 & 564 & $1.0\times10^{-14}$ \\
reduction split into 16 chunks & 0.844 & 948 & $2.9\times10^{-4}$ \\
row sum of the elementwise product & 0.858 & 932 & $1.2\times10^{-7}$ \\
\bottomrule
\end{tabular}
\end{center}

The form we shipped is dominated on both axes at once: it is three times slower
than the third row and no more accurate. The last row is what the solver now
uses in float32. It reaches the same bandwidth as the forward product, it
materializes no intermediate (0.858 ms for 0.80 GB; materializing would cost
about 2.4 ms), it places no divisibility condition on the sample count, and it
is more accurate than the alternative of equal speed by a factor of $2400$,
which matters because the certificate is float64. Float64 keeps the direct form,
being already at $90\%$ of peak.

End to end on the working-set configuration at $\alpha=0.5$ and
$\lambda=0.1\lambda_{\max}$, float32 falls from $0.0564$ s to $0.0467$ s, a
factor of $1.21$, at an unchanged 64 evaluations and an equivalent certified
gap, so the gain is entirely in cost per iteration. Float32 now beats float64
($0.0606$ s) by $1.30\times$, where before it did not, which revises the earlier
row of this section reporting that float32 iterates buy nothing: that row is a
full-width plain FISTA configuration and it stands as measured, but the
conclusion drawn from it does not generalize. The realized $1.21\times$ is
smaller than the $3.07\times$ on the isolated product because the solver runs on
a working set of 125 to 500 columns rather than all 2000. The revised prediction
for a saturated float32 gradient is a factor of two over float64, not the factor
of four predicted above, because two passes over the design are unavoidable.
This is one device. The second machine could not be used to check it, its CUDA
installation being broken pending a reboot.

\paragraph{Full training set (epsilon, $4\times10^5$ samples, Blackwell GPU).}
The training split does not fit next to the language-model server on the
RTX 4090 in float64, so it ran on one RTX PRO 6000 Blackwell GPU with the same
settings ($\alpha=0.5$, $\lambda=0.1\lambda_{\max}$, target gap $10^{-6}$ on our
certificate); CPU baselines on one core. 

\begin{center}
\small
\begin{tabular}{llrrl}
\toprule
Solver & device & seconds & iterations & gap \\
\midrule
cuML (OWL-QN, float32) & Blackwell & 0.755 & 37 & $3.0\times10^{-7}$ \\
ours, FISTA (float64) & Blackwell & 1.53 & 192 & $6.7\times10^{-9}$ \\
ours, forward-backward & Blackwell & 11.2 & 1408 & $7.3\times10^{-7}$ \\
adelie (30-point path) & 1 core & 39.0 & & $5.4\times10^{-8}$ \\
scikit-learn saga & 1 core & 99.2 & & $1.7\times10^{-10}$ \\
skglm (tolerance tightened to $10^{-6}$) & 1 core & 247.9 & & $8.5\times10^{-9}$ \\
\bottomrule
\end{tabular}
\end{center}

On the same GPU cuML is $2.0\times$ faster than our FISTA (against $3.7\times$ on
the test split on the 4090). Our FISTA is $25\times$ faster than the fastest
single-core baseline, which is not a like-for-like hardware comparison. The
logistic certificate is the relative duality gap, which a diverging iterate
cannot pass, so these passes do not share the weakness found in the QP
certificate.

\section{The knowledge-base creation loop}
\label{app:log:kbloop}

 Restarted reflected Halpern PDHG was reached
by transferring general operator ideas (Halpern's anchor, then restarts, then
reflection) to one averaged operator, PDHG. The loop automates that route for
another averaged operator, forward-backward ($\alpha = 2/(4-aL)$ for $a<2/L$),
on elastic-net logistic regression. Ingredients come from the knowledge base:
Krasnoselskii-Mann relaxation, Halpern anchoring, fixed and residual-triggered
restarts, reflection under the anchor, safeguarded Anderson acceleration, and
FISTA momentum as a separate base. Two gate rules were added for this:
residual-triggered restart of the anchor (sound because Halpern's residual
bound forces every segment to end) and Anderson acceleration with the
Zhang, O'Donoghue and Boyd safeguard.

The composer enumerates 97 compositions; the gates reject 58 before anything
runs (for example relaxation beyond $1/\alpha$ without an anchor, momentum
combined with an anchor, two outer extrapolations together) and admit 39,
each run to the certified duality gap. Plain and anchored recipes reproduce
the existing runners exactly (same evaluation counts and iterates), which
validates the composer.

Novelty is checked against the knowledge base's evidence for sparse
generalized-linear-model practice, and the check mostly came back negative:
Anderson acceleration of proximal gradient on logistic regression is
published (Mai and Johansson, ICML 2020), a Halpern anchor has been used on
LASSO (Numerical Algorithms, 2026), and restarting FISTA on such problems is
standard. The only ingredient recorded as unused in that practice, on two
searches, is a periodic restart of a Halpern anchor on forward-backward. A
related June 2026 preprint applies restarted reflected Halpern acceleration
to augmented primal-dual methods (arXiv 2606.16552).

The composer is generic over base operators: every wrapper acts on the
operator's state as a pytree of arrays, so it applies unchanged to forward-backward, FISTA,
PDHG, Condat-Vu, Douglas-Rachford and Davis-Yin, and it reproduces the policy
runner exactly on forward-backward, on Condat-Vu with a restarted anchor for a
real QP, and on PDHG for a transport LP. On the QP instance AUG3DCQP the gates
choose Condat-Vu as the base (19 admitted, 29 rejected) and the loop transfers
Anderson acceleration to it without any QP-specific code: 257 evaluations and
0.030 s against 6400 and 0.179 s for plain Condat-Vu, to a certified $10^{-6}$
KKT residual (both times were measured before the flattening defect of
the timing correction was fixed; the evaluation counts are unaffected). The restarted Halpern chain behind restarted Halpern PDLP loses
there (12,500 to 19,000 evaluations), so that chain is not universally
beneficial. Anderson is recorded as already used in QP practice (SCS uses it),
so this is an automatic rediscovery, not a new algorithm. Family-wide, Condat-Vu + Anderson passes the
certificate on 66 of 137 instances against 55 for plain Condat-Vu; under the
accuracy audit (Section~\ref{app:log:qp}) 55 against 45 are not grossly
wrong; after the composer fixes the rerun gives 64 against 55 passes and 53
against 45 not grossly wrong (Section~\ref{app:log:qp}).

On the epsilon test split at a $10^{-8}$ gap (RTX 4090, after the Anderson accounting fix)
the gates admit 39 of 97 compositions; 33 reach the certificate and 6 hit the cap of
20,000 evaluations, all six unrestarted Halpern anchors, whose $O(1/N)$ rate is too slow
for this accuracy. Plain FISTA is fastest (192 evaluations, 0.947 s). The best
composition on forward-backward uses the step $1.9/L$ (1088 evaluations, 3.90 s, with or
without relaxation); Anderson acceleration lowers the evaluations to 641 to 705 but not
the time (4.8 to 5.0 s); restarted Halpern anchors need 2176 to 2304 evaluations (7.7 to
8.1 s). No composition of these ingredients beats FISTA on this instance: the transfer
route that produced restarted Halpern PDHG finds nothing here.

For the proof-check stage, the performance-estimation package PEPit is
installed in an isolated environment and reproduces the Halpern worst case
exactly ($\|x_N - T x_N\|^2 \le 0.0330579\,\|x_0 - x^\star\|^2$ at $N = 10$); its
fixed-point, monotone-inclusion and composite families are the starting point
for checking a transferred composition before it runs.

That check is now part of the loop .
A fixed-coefficient recipe on an $\alpha$-averaged base, $T = (1-\alpha)I + \alpha N$
with $N$ nonexpansive, relaxed by $\rho$ and iterated plainly or with the
Halpern anchor, becomes a performance-estimation problem whose value is the
worst case of $\|x_N - T x_N\|^2$ over the whole operator class. It reproduces
Halpern's bound $(2/11)^2$ at $N=10$, and it agrees with the gates without being
told what they decided: reflection without an anchor ($\rho\alpha = 1$), which the
gate rejects, shows no worst-case decay from $N=10$ to $N=20$, while strict
relaxation and reflection under the anchor, which the gate admits, contract.
The gate and this analysis are two independent derivations of the same
admissibility, one symbolic and one numerical. Restarts and Anderson
acceleration have iterate-dependent coefficients, so they are not
fixed-coefficient problems and are left to the certified runs.

\section{Diagnose, mutate, re-run: closing the gap to cuML}
\label{app:log:diagnose}

The composition loop of Section~\ref{app:log:kbloop} found nothing that beats
FISTA on epsilon, so a diagnosis stage was added in front of it. It splits the
gap to the baseline into the two factors of $T = N \times t$ and turns
measurements into symptoms, each mapped to knowledge-base ingredients that act
on one factor:
\begin{itemize}
\item the solution is sparse (197 of 2000 features) but every step reads all
 columns: Gap Safe screening with mid-solve re-planning (acts on $t$);
\item the step is memory bound (about 90\% of peak bandwidth in float64):
 float32 storage, or a fused gradient kernel where the float32 transposed
 product is slow, as on the 4090 but not on the Blackwell (acts on $t$);
\item the baseline needs far fewer iterations (37 against 192): variable
 metric, Anderson acceleration, a quasi-Newton metric (acts on $N$). Anderson
 and the Halpern anchor are now reachable, by routing the configuration to the
 composition engine of Section~\ref{app:log:kbloop}, and both were tried and
 beaten on measurement; the quasi-Newton metric and the fused kernel remain
 unimplemented.
\end{itemize}
Further rules compare a run with the run it was mutated from: a re-plan that
restarts FISTA's momentum (keep it); exact-size re-plans that compile a new
shape inside the solve (pad to halving buckets compiled from $(n, p)$ alone);
iteration counts rounded up to the certificate cadence (halve $K$); safe
screening that cannot act until the gap is small (start from a working set
ranked by dual scores and let the full-problem certificate add violators, the
approach of celer and skglm); and, on a working set, the exact Lipschitz
constant of the selected columns in place of power iteration. The search is a
beam of width two over single mutations, so a step that pays only after the
next one (screening, before its re-planning costs are fixed) is not discarded;
a first greedy version stopped at 0.500 s for exactly that reason. The
reported chain is the ancestry of the best configuration. Every certificate is
the full-problem duality gap, so no mutation changes what is certified, only how
fast it is reached.

Epsilon test split on the RTX 4090, target gap $10^{-6}$, medians of five
solves after an untimed warm-up, step size timed:

\begin{center}
\small
\begin{tabular}{lrrl}
\toprule
Configuration & seconds & evaluations & gap \\
\midrule
FISTA, $K=64$ & 0.775 & 192 & $6.3\times10^{-9}$ \\
FISTA, $K=16$ & 0.642 & 144 & $9.0\times10^{-7}$ \\
\quad + Gap Safe screening, momentum kept, buckets & 0.456 & 96 & $5.5\times10^{-7}$ \\
\quad + working set from $p/16$ columns & 0.081 & 48 & $1.0\times10^{-9}$ \\
\quad + certificate every 32 steps, exact step on the working set & 0.067 & 64 & $9.4\times10^{-7}$ \\
cuML (OWL-QN, float32), same GPU, eight runs & 0.166 to 0.295 & 35 to 39 & $1.6$ to $4.6\times10^{-7}$ \\
\bottomrule
\end{tabular}
\end{center}

The median of the last five cuML runs on the 4090 is 0.227 s; one of them
needed a tolerance of $10^{-7}$ to reach the target. On the Blackwell training
split (one RTX PRO 6000 GPU; cuML alone on an identical GPU, and five times on
the same GPU) the rows read: FISTA 1.775 s, screening 1.026 s (96 evaluations),
working set 0.176 s (48 evaluations, gap $9.5\times10^{-10}$), and 0.142 s with
the certificate every 32 steps and the exact step. cuML took 0.482 to 0.636 s
after a cold first run of 1.32 s (median 0.615 s, 34 to 43 iterations), and two
of the five runs needed a tolerance tighter than $10^{-4}$ to reach the target.
So with nothing excluded but compilation, the best configuration found by hand
is $2.5\times$ faster than cuML's fastest run on the 4090 and $3.4\times$ on the
Blackwell ($3.4\times$ and $4.3\times$ against cuML's medians); the loop's own best
configurations, below, are $3.0\times$ and $3.7\times$ faster than cuML's fastest runs
($4.1\times$ and $4.7\times$ against the medians).

What the loop found by itself: on the 4090, in its first version (with $L$
still computed outside the timer), the beam search from plain FISTA at $K=64$
reached $K=32$ with screening, a working set, buckets, float32 storage and
kept momentum, 0.127 s against cuML's 0.166 s. On the Blackwell, with the step
size timed, it reached the working set in its third round (0.333 s, then 0.232 s
at $K=32$, against 0.571 s). With the step size timed, the 4090 loop took screening, the working set,
$K=32$ and float32 storage, 0.102 s ($1.6\times$ cuML's fastest run), and the
Blackwell loop stopped at 0.232 s; both fell short of the best configuration
found by hand. The shortfall was traced to a gene no rule could move, the size
of the first working set, which held 250 columns for a solution with 197
nonzeros. A rule was added that halves it when it already holds the final
support. The Blackwell rerun then went, unaided, from 1.777 s through screening
(1.518 s), the working set (0.333 s), the smaller first set (0.191 s), $K=32$
(0.141 s) and float32 storage (0.132 s): ahead of the configuration found by
hand, and $3.7\times$ faster than cuML's fastest run on the same GPU model. The
4090 rerun with the new rule went from 0.777 s through screening (0.713 s), the
working set (0.148 s), the smaller first set (0.086 s), $K=32$ (0.063 s), float32
storage (0.059 s) and kept momentum (0.056 s), also ahead of the configuration
found by hand (0.067 s): $3.0\times$ faster than cuML's fastest run on the same GPU
and $4.1\times$ faster than its median. This is the harness evolving from
its own failures, the pattern CAKE (arXiv 2608.12629) applies to kernel agents,
where recurring failures become verifier rules.

\paragraph{Guided against blind mutation (Blackwell, epsilon training split).}
 Both arms search the same eight genes
(certificate cadence $K$, screening, momentum on re-planning, buckets, working
set, size of the first working set, step rule, precision) from the same start
(FISTA, $K=64$, power-iteration step), with the same beam of two, the same
certificate and the same budget of 24 evaluated configurations. The guided arm
takes its mutations from the diagnosis rules; the blind arm takes four random
single-gene changes per expanded configuration and sees only certified wall
time. Three seeds, one per GPU; medians [min, max] of the best certified time
(cuML 0.571 s): after 8 configurations, guided 0.802 s [0.801, 0.805] and blind
1.207 s [0.910, 1.239]; after 16, 0.190 s [0.186, 0.192] and 0.192 s [0.140,
0.782]; after 24, 0.140 s [0.136, 0.145] and 0.141 s [0.140, 0.248].

On this small space the diagnosis does not reach a better endpoint: one gene
(the working set) carries most of the gain, and random single-gene changes find
it within 24 configurations in two of three seeds. What the diagnosis buys here
is reliability and early progress: its three runs agree to within 7\%, while the
blind runs span 0.14 to 0.78 s at 16 configurations and one ends at 0.248 s. The
stronger evidence for the diagnosis lies before this ablation: the timing-only
composition loop of Section~\ref{app:log:kbloop} searched a space without
screening or working sets and found nothing that beats FISTA, and those
ingredients entered the search space only because the diagnosis named the
sparse solution and the gap-limited screening as the bottleneck. On the RTX 4090
(epsilon test split, cuML 0.166 s, three seeds run in sequence) the pattern
repeats: after 16 configurations guided 0.087 s [0.086, 0.087] and blind
0.150 s [0.065, 0.469]; after 24, guided 0.064 s in all three seeds and blind
0.065 s [0.064, 0.119]. On both devices the diagnosis gives the same endpoint
as a lucky blind run, every time; the blind search gives it in two runs of three.

Does the diagnosis matter more when the space a blind search must cover is
larger? The blind arm was rerun on the 4090 with a larger space of real solver
settings (six cadences and six first-set sizes instead of four, plus the number
of power iterations, the screening threshold, and FISTA against
forward-backward); the guided arm is unchanged, because its rules never propose
those genes. Three seeds, best certified time after 16 configurations: blind
0.149 s [0.135, 0.651] against 0.150 s on the small space; after 24, blind 0.087 s
[0.063, 0.119] against 0.065 s on the small space and 0.064 s for the guided arm.
The blind median gets worse as the space grows ($1.36\times$ the guided time after
24 configurations) while one blind seed still finds the best configuration. With
three seeds this is a trend, not a firm result.

Three more seeds on the 4090 (six in all, each running all three arms) firm up
one conclusion and soften another. After 24 configurations the medians are
guided 0.065 s [0.064, 0.066], blind 0.077 s [0.064, 0.119] on the small space and
0.080 s [0.059, 0.119] on the large one; after 16, 0.086 s [0.086, 0.093], 0.099 s
[0.065, 0.469] and 0.142 s [0.078, 0.651]. The guided arm comes within 5\% of
0.064 s in 6 of 6 seeds, each blind arm in 2 of 6: the diagnosis makes the search
reliable. The larger space hurts the blind search clearly after 16
configurations and only slightly after 24, so the effect of space size is not
established. One blind seed on the large space beat every guided run: 0.059 s in
32 evaluations, from a finer cadence ($K=8$) and a cheaper step estimate (10 power
iterations instead of 30) on the working set. No rule proposed either change: the
cadence rule stopped at $K=32$ and no rule touched the number of power
iterations. Two rules were added for these, again a failure of the loop turned
into a rule, and the guided arm was rerun (six seeds each): with the rules as
first written (a condition error kept the cadence rule from ever firing), fixed,
and fixed with the guided arm limited to four proposals per expanded
configuration, rules for the dominant factor of $T=N\times t$ first. None moved the
median after 24 configurations (0.0660, 0.0648 and 0.0648 s against 0.0647 s), and
the 0.059 s configuration was reached at most once. A direct timing (ten
interleaved solves per configuration) shows why and settles whether the blind
result was noise. The guided best takes 0.0667 s [0.0633, 0.0718]; the blind best
0.0568 s [0.0528, 0.0686], a real gain; the guided best with only the cheaper step
estimate 0.0492 s [0.0468, 0.0534], the fastest configuration found so far
($3.4\times$ cuML's fastest run on this GPU); and the guided best with only the finer
cadence 0.0891 s [0.0821, 0.0978], worse. So the step-estimate rule is right and
points directly at the best configuration, but within 24 evaluations the guided
search spends its budget on the other proposals before it reaches it, while the
cadence rule is wrong and has been removed. Raising the budget to 32 configurations (six seeds,
step-estimate rule kept, cadence rule removed) leaves the median at 0.0649 s
[0.0629, 0.0662] and reaches the fastest configuration in no seed, and the records
show exactly why: on every seed the search first reaches the configuration it
must extend (the working set started at $p/16$ with $K=32$) at evaluation 21, and
expands it only in the last round, where the rules propose keeping momentum
(evaluation 30), a finer cadence (evaluation 31) and then the cheaper step
estimate, one evaluation past the budget. The search is deterministic apart from
timing noise, so the miss is structural: the rule is right, the order in which a
node's proposals are tried is not. Each new rule makes a fixed budget
thinner; ranking proposals by a calibrated prediction of their gain, the role of
CAKE's cost model, is the missing piece. Neither rule was used in the
generalization stage below, whose tuning procedure was fixed before they existed.

Working sets with dual-score selection are published practice (celer,
skglm, Blitz), so that part is an automatic rediscovery of known structure
exploitation, not a new algorithm; what is new is that the chain was derived
from measurements, per device, with the certificate unchanged. The ablation
below separates the ingredients, and the one carrying the speedup is not the
working set itself but the step taken from the Lipschitz constant of the
current working set, which that literature does not do: it is almost entirely
coordinate descent, where the step is per feature and computed once on the
whole matrix, so a restricted constant changes nothing there. Ours runs in
float64, cuML in float32. Compilation (2.7 to 3.9 s) depends only on $(n, p)$
and is excluded here, as is cuML's warm-up solve. That exclusion was misleading
when these rows were measured, because the programs were rebuilt on every call
and so the cost was paid on every solve; they are reused across
calls and across penalty weights, so a repeated solve and a path both pay it
once, and a first solve still pays it in full. cuML's time varied between runs of the same solve,
and the comparison uses its fastest. At looser tolerances cuML does not reach
the target (gap 0.32 at $10^{-2}$, $9\times10^{-3}$ at $10^{-3}$), so $10^{-4}$ is its
fastest passing setting. These rows are one instance at one penalty weight; the
ablation and the held-out stage below extend the comparison to six and ten.

\paragraph{Baseline precision.}
Our worker first ran cuML in float32 only (it cast the data before fitting). On
harder instances float32 cannot reach the target: on the epsilon test split at
$\alpha=0.9$, $\lambda=0.02\lambda_{\max}$ its gap stays at $1.4\times10^{-5}$ or
above for every tolerance from $10^{-4}$ to $10^{-8}$, and at $\alpha=0.5$,
$\lambda=0.02\lambda_{\max}$ at $1.4\times10^{-6}$ or above, while cuML in float64
certifies both (1.414 s, 317 iterations; 0.567 s, 127 iterations). On the
instance of the table above both precisions pass, float64 in 0.211 s in one
run, slower than float32's fastest (0.166 s), so the comparisons above stand. The
worker now takes the precision as an option, and every later comparison uses,
per instance, the faster of the two precisions that reaches the target. Counting
an instance where float32 fails as a win would reward our float64 arithmetic,
not our method.

\paragraph{The loop run end to end.}

Starting from
plain forward-backward with Nesterov momentum on the epsilon test split at
$\alpha=0.5$, $\lambda=0.1\lambda_{\max}$, which loses to cuML by $4.8\times$
($0.770$ s against $0.159$ s), the search reached $0.059$ s in four diagnosed
steps: screening ($0.590$ s), a working set ($0.144$ s), a smaller initial
working set ($0.081$ s), and a halved certificate cadence ($0.059$ s). That is
$2.67\times$ faster than the baseline it started behind. Each step is one
symptom mapped to one typed mutation, and the certificate is the full-problem
duality gap throughout, so nothing in the chain changes what is being certified.
The operator ingredients were proposed, admitted by the gate, run, and rejected
on measurement rather than by assumption: the anchored iteration could not
certify within twenty times the incumbent's iteration count and was stopped
there, and safeguarded Anderson was slower by a factor of fifty. On this cell
the gain is entirely on the screening axis.

Bounding the search was necessary and is worth recording, because two earlier
attempts at the same run exhausted their time budget. A step that loses badly is
followed once and then neither re-proposed from other parents nor expanded into
its own descendants, and a composed recipe is cut off once it exceeds twenty
times the best certified iteration count, which is a competitiveness bound
rather than a convergence claim. An earlier attempt also stopped one step short,
at $1.93\times$, for a reason that was our instrument and not the method: several
candidates in the last round failed for want of device memory, and those
failures were recorded as slow configurations rather than as failures to
measure. Distinguishing the two did not merely correct the record, it let the
search find the cadence step and a configuration $1.4\times$ faster than the one
it had settled for.

\section{Isolating the ingredients}
\label{app:log:ablation}

 The chain above mixes several changes, so
this stage varies one at a time. Ten instances of the epsilon test split,
$\alpha \in \{0.5, 0.9\}$ and $\lambda/\lambda_{\max} \in \{0.02, 0.05, 0.1, 0.2,
0.5\}$, five arms, seven repetitions each. The repetitions are interleaved, so
that repetition $k$ of every arm runs next to repetition $k$ of the others:
pairing arms by repetition index means nothing when each arm's repetitions run
in one block minutes apart, which is how an earlier version of this experiment
was arranged, and the paired intervals below are reported only because the
interleaved design earns them.

The arms are plain forward-backward with momentum; a working set whose step
comes from the Lipschitz constant of the whole matrix, with momentum either
reset or carried across changes of the set; and a working set whose step comes
from the Lipschitz constant of the current set, again with momentum reset or
carried. Three times are reported separately throughout, because they are
three different claims: the solve loop, a repeated solve (the loop plus
compilation, with compiled programs reused), and a first solve on an empty
cache.

\begin{center}
\small
\begin{tabular}{lcc}
\toprule
Contrast & geometric mean & $95\%$ interval \\
\midrule
reduced-set step against whole-matrix step, momentum reset & $2.058\times$ & $[1.635, 2.590]$ \\
reduced-set step against whole-matrix step, momentum carried & $1.840\times$ & $[1.426, 2.374]$ \\
momentum carried against reset, whole-matrix step & $1.171\times$ & $[1.103, 1.243]$ \\
momentum carried against reset, reduced-set step & $1.047\times$ & $[0.992, 1.105]$ \\
both against neither & $2.154\times$ & $[1.656, 2.801]$ \\
both against plain forward-backward & $11.687\times$ & $[8.227, 16.601]$ \\
\bottomrule
\end{tabular}
\end{center}

The step taken from the Lipschitz constant of the current working set is what
carries the result: its interval excludes one on every instance and across
them. Carrying momentum is real only while the step comes from the whole
matrix; once the reduced-set step is in place the interval straddles one and
the within-instance comparison is inconclusive on seven of the ten. The reason
is visible in the counts: the reduced-set step cuts a solve from 288
evaluations and three re-plans to 96 and two, and momentum can only be worth
what the re-plans would otherwise destroy. Against cuML, taking per instance
the faster of its two precisions that reaches the target and its fastest run of
that precision, the best arm is $2.28\times$ with a $95\%$ interval of $[1.73,
3.01]$ and faster on ten of ten for the solve loop and for a repeated solve,
and $0.06\times$ and faster on none for a first solve. Running every arm in
float32, which matches the precision cuML wins with on most of these
instances, gives $2.43\times$ and does not change any of the contrasts.

One negative is worth stating because it was the hypothesis this grid was run
to test. The advantage over cuML grows with the penalty weight, so it would be
convenient if the ingredient that carries the speedup grew with it too. It does
not. At $\alpha = 0.5$ the gain from the reduced-set step falls from
$1.96\times$ at $0.02\lambda_{\max}$ to $1.13\times$ at $0.5\lambda_{\max}$,
while the advantage over the baseline rises across that same range, so the two
move in opposite directions; at $\alpha = 0.9$ the gain is not monotone at all
and peaks in the middle. The two effects have different explanations. Heavily
penalized instances are already short, thirty-two evaluations at
$0.5\lambda_{\max}$, so a larger step has little left to save and the advantage
there comes from the working set being small, which is a cost-per-iteration
effect; lightly penalized instances give the larger step the most room in
relative terms while remaining hard in absolute terms, which is where the
baseline closes. The ingredient acts on the iteration count and the regime
boundary is set by the cost per iteration.

\section{Generalization: a frozen portfolio on held-out instances (not critical)}
\label{app:log:general}

 Following
the separation CAKE (arXiv 2608.12629) makes between tuning one shape and building a
library, the domain and its split were written to before any
tuning: four real LIBSVM datasets (epsilon training split,
$4\times10^5\times2000$; epsilon test split, $10^5\times2000$; gisette, $6000\times5000$;
covtype, $581012\times54$), $\lambda/\lambda_{\max}\in\{0.5, 0.2, 0.1, 0.05, 0.02\}$ and
$\alpha\in\{0.5, 0.9\}$. The 18 tuning instances are epsilon training, gisette and
covtype at 0.5, 0.1 and 0.02; the 22 held-out instances are the same datasets at 0.2
and 0.05 (unseen regularization) and every epsilon test instance (an unseen shard of a
tuned source). The guided loop ran on every tuning instance on the Blackwell, step
size timed, including the 9 on which cuML in float32 never reaches the target (there
it is tuned against cuML's closest run, which only feeds the rules). Every distinct
best configuration was then timed on every tuning instance, and a portfolio of at most
three recipes with a dispatch tree of depth two on declared features ($\log n$,
$\log p$, $\lambda/\lambda_{\max}$, $\alpha$) was fitted to the geometric-mean time on the
tuning instances only, then frozen with a content hash (\texttt{ae2a840e\ldots}) before
any held-out instance ran.

The fitted portfolio has two recipes and one split, on $\lambda/\lambda_{\max}\le0.06$:
below it, a working set started at $p/64$ with screening, buckets and $K=64$; above it,
the same with $K=16$ and float32 storage. On the tuning instances its geometric-mean
time is 0.071 s, against 0.084 s for the single best recipe and 0.060 s for an oracle
that picks the best candidate per instance.

Held-out result against cuML in both precisions (three runs per instance, fastest
shown; our time is the median of three solves after an untimed warm-up; the last
column divides the faster cuML precision that reaches the target by our time):

\begin{center}
\small
\begin{tabular}{lllrrrr}
\toprule
dataset & $\alpha$ & $\lambda/\lambda_{\max}$ & ours (s) & cuML float32 (s) & cuML float64 (s) & speedup \\
\midrule
covtype & 0.5 & 0.05 & 0.1102 & fails ($\ge4\times10^{-6}$) & 0.3258 & 2.96 \\
covtype & 0.5 & 0.2 & 0.0552 & 0.0698 & 0.0778 & 1.26 \\
covtype & 0.9 & 0.05 & 0.2265 & fails ($\ge3\times10^{-5}$) & 0.6153 & 2.72 \\
covtype & 0.9 & 0.2 & 0.0670 & fails ($\ge2\times10^{-5}$) & 0.1330 & 1.98 \\
epsilon test & 0.5 & 0.02 & 0.2380 & 0.3296 & 0.3679 & 1.39 \\
epsilon test & 0.5 & 0.05 & 0.1806 & 0.2085 & 0.2931 & 1.15 \\
epsilon test & 0.5 & 0.1 & 0.0499 & 0.1296 & 0.1587 & 2.60 \\
epsilon test & 0.5 & 0.2 & 0.0388 & 0.1372 & 0.0799 & 2.06 \\
epsilon test & 0.5 & 0.5 & 0.0103 & 0.1465 & 0.0638 & 6.18 \\
epsilon test & 0.9 & 0.02 & 0.2730 & fails ($\ge5\times10^{-6}$) & 0.8840 & 3.24 \\
epsilon test & 0.9 & 0.05 & 0.1286 & fails ($\ge1\times10^{-6}$) & 0.4922 & 3.83 \\
epsilon test & 0.9 & 0.1 & 0.0697 & 0.3237 & 0.3375 & 4.65 \\
epsilon test & 0.9 & 0.2 & 0.0327 & 0.1879 & 0.2496 & 5.74 \\
epsilon test & 0.9 & 0.5 & 0.0130 & 0.1180 & 0.0786 & 6.03 \\
epsilon train & 0.5 & 0.05 & 0.6301 & 0.6861 & 0.8139 & 1.09 \\
epsilon train & 0.5 & 0.2 & 0.1877 & 0.4417 & 0.1857 & 0.99 \\
epsilon train & 0.9 & 0.05 & 0.4403 & fails ($\ge3\times10^{-6}$) & 1.7208 & 3.91 \\
epsilon train & 0.9 & 0.2 & 0.1326 & 0.6295 & 0.6359 & 4.75 \\
gisette & 0.5 & 0.05 & 0.0313 & 0.1160 & 0.2050 & 3.70 \\
gisette & 0.5 & 0.2 & 0.0194 & 0.0585 & 0.1072 & 3.02 \\
gisette & 0.9 & 0.05 & 0.0443 & fails ($\ge3\times10^{-6}$) & 0.3383 & 7.64 \\
gisette & 0.9 & 0.2 & 0.0178 & 0.0800 & 0.1352 & 4.49 \\
\bottomrule
\end{tabular}
\end{center}

Our dispatched recipe certifies all 22 held-out instances, and cuML in float64
reaches the target on all 22 as well, so every instance has a baseline. Against
the faster cuML precision that reaches the target on each instance, the frozen
portfolio is faster on 21 of 22, by a geometric mean of $2.91\times$ against cuML's
fastest run and $3.13\times$ against its median (range $0.99\times$ to $7.64\times$;
the one loss, epsilon training at $\alpha=0.5$, $\lambda=0.2\lambda_{\max}$, is 0.1877 s
against 0.1857 s). Against float32 alone, on the 15 instances where it reaches the
target, the figures are $3.21\times$ and $3.82\times$ (15 of 15); against float64 alone,
$3.33\times$ and $3.42\times$ (21 of 22). For reference, the single best recipe gives
$2.37\times$ (20 of 22) and plain FISTA $0.36\times$ (3 of 22) against the faster cuML
precision, so the gain comes from the recipes the loop found and from dispatching
between two of them, not from FISTA. The tuning used the rules as they stood before
the two rules of Section~\ref{app:log:diagnose} were added; the cheaper step estimate
those runs identified is not in this portfolio.

This comparison covers one regularization family (elastic-net logistic), four dense datasets and
one GPU type; cuML's times vary between runs (the table shows the fastest of three);
ours runs in float64 or float32 storage with a float64
certificate, and cuML was given both precisions.

The exclusion of compilation above is not the symmetric statement it appears to be,
and the numbers in this subsection must be read with the following correction. Our
compilation is about three seconds against solves of $0.01$ to $0.4$ s, and it was re-paid on every call, because the compiled programs were held in a
dictionary local to the solve and discarded when it returned. The excluded cuML
warm-up is a two-iteration fit on a small slice. So the figure reported above was the
solve loop alone, and no caller could obtain it: every solve, including every point of
a regularization path, paid the compilation again. Programs are now reused across
calls, keyed on a fingerprint of the labels as well as the shapes, since the labels and
the two penalty parameters are constants in the compiled program. A repeated solve at
one penalty weight now compiles once, and so does a path over several weights: the two
penalty parameters were made arguments of the compiled program rather than constants
in it. Solving a five-point path now costs one compilation instead of five, and the
solutions are unchanged to the last bit.

Three times are therefore distinct and they are three different claims: a first solve
on an empty cache, a repeated solve, and the solve loop. Re-measured under that
accounting, on the ten held-out instances that can still be measured, the dispatched
portfolio is $2.39\times$ the better cuML precision's fastest run, with a $95\%$
interval over the ten instances of $1.40$ to $4.07$, and faster on eight of ten; the
repeated solve is identical, and these records do not measure a first solve, because
the warm-up absorbs the compilation. Variation between runs is small next to variation
between instances, which is what makes an interval over instances the right summary
here: our three repetitions span at most $1.06\times$ and cuML's at most $1.57\times$,
and the comparison takes cuML's fastest run rather than its median, which favours the
baseline. It is an interval over instances and not a paired test between the two
solvers, whose repetitions are independent runs with no matched pairing. The advantage grows monotonically with
the penalty weight, from $1.50\times$ at $\lambda = 0.02\lambda_{\max}$ to
$7.00\times$ at $0.5\lambda_{\max}$, and both losses fall at the weakest penalties,
where the solution is densest and screening removes least. The eight-fold spread is
visible in the iteration counts, which run from 16 at $0.5\lambda_{\max}$ to 448 at
$0.02\lambda_{\max}$. The claim this supports is therefore about well-regularized
instances of the family, with a crossover near $0.05\lambda_{\max}$ at $\alpha=0.5$,
not about the family as a whole.

Only ten of the twenty-two held-out instances could be re-measured. Two of the four
datasets are no longer on disk and are absent from the data manifest, and the largest
does not fit in this GPU's memory beside the language-model server, needing
$5.98$ GiB for the design matrix alone. The earlier $2.91\times$ over twenty-two
instances and the $2.39\times$ here are not the same quantity measured twice.

\section{Variable metric ideas considered for QP}
\label{app:log:varmetric}

Source: J.F.~Bonnans, J.Ch.~Gilbert, C.~Lemar\'echal and C.A.~Sagastiz\'abal, A
family of variable metric proximal methods, Mathematical Programming 68
(1995) 15--47, read in full. The first idea is now implemented and measured
(Section~\ref{app:log:adaptive}). The quasi-Newton primal metric is implemented as
the builder \texttt{condat\_vu\_metric} and measured on all 137 QP instances: it does not
help. With ten Krylov secant pairs it solves 55 instances like the scalar step (one
gained, one lost) and needs $1.20\times$ the iterations and $1.52\times$ the time where
both solve; built from warm-up iterates it solves 53 (two gained, four lost), with
$1.16\times$ the iterations. The remaining ideas are recorded with the gate rule they
would need.

\begin{itemize}
\item A gate for adapting the metric. The paper's global result allows any
 sequence of positive definite metrics $M_n$ in the proximal step, provided a
 sufficient-decrease test holds and $\sum_n \lambda_{\min}(M_n^{-1}) = \infty$
 (for BFGS metrics this follows from $\operatorname{tr} M_n \le (n+1)C$). For
 the splitting schemes in our genome the matching condition is the variable
 metric quasi-Fej\'er rule of Combettes and V\~u (2013, 2014): uniform bounds
 $\alpha I \preceq M_n \preceq \beta I$ and $M_{n+1} \preceq (1+\eta_n) M_n$
 with $\sum_n \eta_n < \infty$. This is the missing certificate for the
 adaptive step ratio: PDLP's primal weight and OSQP's adaptive $\rho$ are
 scalar varying metrics adapted without such a budget.
\item A quasi-Newton primal metric. With secant pairs $s_n = x_{n+1}-x_n$ and
 $y_n = \nabla f(x_{n+1}) - \nabla f(x_n)$, which for a QP equal $P s_n$
 exactly, a limited-memory BFGS metric can drive the primal step. The
 structural fact that makes this cheap: in the OSQP form of a QP the
 separable primal term is zero, so the primal update
 $x^+ = x - M_n^{-1}(\nabla f(x) + A^\top u)$ needs no proximal map in the
 non-diagonal metric, only an application of $M_n^{-1}$. The step condition
 in the new metric must be re-checked (power iteration on $M_n^{-1}P$ and
 $A M_n^{-1} A^\top$), which is affordable only at chunk boundaries: metric
 updates would happen at the certificate cadence $K$, the one gene already
 known to couple $N$ and $t$.
\item Reflection with a matched metric. When the metric equals the Hessian of
 a quadratic, the proximal step moves only half way to the minimizer, and the
 step $x_{n+1} = x_n + 2(x^p_n - x_n)$ recovers it; the paper's line search
 prefers this step and falls back when a descent test fails. In our genome
 this is relaxation $\rho=2$ combined with the safeguarded switch wrapper.
 Reflected Halpern restarts without a matched metric did not help on
 transport (Section~\ref{app:log:ot}).
\item A relative acceptance test for inexact proximal steps. An approximate
 proximal point is accepted when a lower model of the objective is accurate
 enough there, a relative test that needs no summable error sequence. This
 is the stopping rule the conjugate-gradient option of a factorization-based
 proximal step would need, and it connects to the relative-error criterion
 of Solodov and Svaiter in the knowledge base. It is not new for QP:
 rPDHCG solves the primal subproblems of PDHG inexactly by conjugate
 gradients, and a 2026 GPU conic-QP solver uses a relative proximal-error
 budget for restarted PDHG (arXiv 2608.09159).
\item Prior use, from one search each: a quasi-Newton primal metric appears
 in general proximal splitting (arXiv 1801.08691, 2209.14019) but was not
 found in LP or QP PDHG practice, so it is a frontier candidate pending a
 wider search; reflection is already used (reflected Halpern PDHG for LP,
 arXiv 2407.16144).
\end{itemize}

\section{A dense row in the wrong place: diagnose and relocate, September 14}
\label{app:relocate}

Portfolio was the one QP family we lost, at every size and every setting we
tried. The cause is geometric rather than numerical. After Ruiz
equilibration the constraint operator of the OSQP portfolio form is
effectively rank one at the top: the ratio of its two largest singular
values is $16$ at $n_{\mathrm{var}}=5{,}070$ and grows to $41$ at
$550{,}000$, and a single row, the budget row $\mathbf{1}^\top x = 1$,
carries $99.9\%$ of the top left singular direction. That row is dense
(density $0.986$) while every other row of the operator has density between
$0.0005$ and $0.09$.

\paragraph{Rescaling the row is not a substitute.}
The obvious objection is that a dense row is exactly what a diagonal
preconditioner is for. It is not, and the distinction is the point of this
subsection. Scaling row $j$ of $A$ by $\varepsilon$, together with $l_j$ and
$u_j$, gives an equivalent problem with a smaller $\|L\|$, so the norm can
be made as small as one likes. What that buys is another matter. On
portfolio at $n=5{,}000$, judged throughout against the original problem:

\begin{center}
\begin{tabular}{lrrrr}
\toprule
$\varepsilon$ & $\|L\|$ raw & iters raw & $\|L\|$ after Ruiz & iters after Ruiz \\
\midrule
$1$    & 72.13 & $>$200{,}000 & 34.65 & $>$200{,}000 \\
$0.3$  & 21.66 & $>$200{,}000 & 34.61 & $>$200{,}000 \\
$0.1$  & 7.29  & 137{,}000    & 34.57 & $>$200{,}000 \\
$0.03$ & 5.74  & 101{,}500    & 34.53 & $>$200{,}000 \\
$0.01$ & 5.74  & $>$200{,}000 & 34.49 & $>$200{,}000 \\
\bottomrule
\end{tabular}
\end{center}

The same sweep on the SVM dual at $n_{\mathrm{var}}=5{,}500$ has the same
shape: $>$200{,}000 iterations at $\varepsilon=1$, a best of $13{,}000$ at
$\varepsilon=0.03$, $20{,}250$ at $\varepsilon=0.01$, and $\|L\|$ after Ruiz
pinned between $70.35$ and $70.67$ with no value converging.

Three things follow, on both families. The scaling does shrink $\|L\|$, by
$12.6\times$ here, and it does help, but the best hand-tuned value never
matches relocation: $101{,}500$ iterations against at most $250$ on
portfolio and $13{,}000$ against $2{,}250$ on the SVM dual, all four figures
measured with plain Condat-Vu so the comparison is like for like. The effect
is not monotone, and the optimum sits at $\varepsilon=0.03$ on both
families, so any rule of the form ``scale the dense row down'' would have to
search for the constant rather than derive it. And Ruiz equilibration, the
preconditioner actually in use, is invariant to the whole sweep, holding
$\|L\|$ between $34.49$ and $34.65$ on portfolio across a hundredfold change
in $\varepsilon$, and never converging on either family. The margin between
the best achievable scaling and relocation is itself family-dependent, a
factor of at least $406$ on portfolio against $5.8$ on the SVM dual, so the
claim is that rescaling cannot replace relocation, not that it is uniformly
hopeless.

The reason is that the dual step $s = 0.49/(t\|L\|^2)$ is a single global
scalar. The dense budget row and the sparse factor rows want different dual
step sizes, and one scalar step with a diagonal scaling cannot serve both;
shrinking $\|L\|$ far enough to help the budget row destabilizes the factor
rows, which is the $\varepsilon=0.01$ line. Relocation does not rebalance
that trade-off, it removes it, by taking the row out of the dual problem
altogether. The difference is categorical, not quantitative.

\paragraph{The transform, derived rather than assumed.}
Given the row the diagnosis selects, the harness reads the transform off
$(A,l,u)$: the row must be an equality, its support must be a prefix block,
every variable it touches must carry simple bounds supplied by identity
rows, and the remaining rows must be equalities. Then
$\{lo \le x \le hi,\ a^\top x = b\}$ is exactly a box-hyperplane projection,
and the special case $a=\mathbf{1}$, $b=1$, $lo=0$, $hi=\infty$ is the
probability simplex with a cheaper exact projection. The procedure returns
``not relocatable'' when any condition fails, which is what keeps it a rule
rather than a table of families.

\paragraph{Two families.}
On portfolio the speedup over Clarabel at tolerance $10^{-6}$ is monotone in
size, and the family is lost below a crossover near $n_{\mathrm{var}} =
100{,}000$:

\begin{center}
\begin{tabular}{lrr}
\toprule
$n_{\mathrm{var}}$ & speedup & $95\%$ interval over 3 seeds \\
\midrule
5{,}070   & $0.022\times$ & \\
50{,}223  & $0.22\times$  & $[0.19, 0.26]$ \\
100{,}316 & $0.90\times$  & $[0.55, 1.45]$ \\
250{,}500 & $1.89\times$  & $[1.52, 2.36]$ \\
500{,}707 & $3.09\times$  & $[2.25, 4.24]$ \\
\bottomrule
\end{tabular}
\end{center}

The second family is the dual of the soft-margin support vector machine
with intercept, whose feasible set $\{0 \le \alpha \le C,\ d^\top \alpha =
0\}$ is a box-hyperplane. It is a different problem from the primal SVM
family used elsewhere in this paper, which has no intercept and therefore no
equality row at all; the two were checked against each other by strong
duality, with primal and dual optima summing to zero within $4\times
10^{-11}$ at three sizes. The diagnosis fires on it with no family-specific
code (row mass $0.9999$, singular gap $35.4$), derives the box-hyperplane
form, and wins at every size we ran, by a margin that grows:

\begin{center}
\begin{tabular}{lrrl}
\toprule
$n_{\mathrm{var}}$ & speedup & $95\%$ interval & \\
\midrule
38{,}500  & $0.79\times$ & $[0.72, 0.87]$ & loses \\
55{,}000  & $1.45\times$ & $[1.26, 1.67]$ & \\
88{,}000  & $2.45\times$ & $[2.24, 2.67]$ & \\
220{,}000 & $5.29\times$ & $[5.09, 5.50]$ & \\
550{,}000 & \multicolumn{2}{l}{measurement in progress} & \\
\bottomrule
\end{tabular}
\end{center}

All figures are three seeds under the accounting of
Section~\ref{app:accounting}, which charges every rung of the tolerance
ladder. The family crosses $1.0$ between $n_{\mathrm{var}} = 38{,}500$ and
$55{,}000$, well below portfolio's crossover, and the margin grows
monotonically above it. Accuracy is better than the baseline on the
objective and comparable on feasibility: at $220{,}000$ our objective gap is
$1.2\times10^{-11}$ against Clarabel's $9.5\times10^{-8}$, with row
violations $3.6\times10^{-7}$ and $1.3\times10^{-7}$ respectively.

\paragraph{Setup dominance reverses at scale.}
The claim elsewhere in this paper that setup rather than iteration count
dominates our runtime is size-dependent, and this family shows where it
turns over. At $n_{\mathrm{var}} = 55{,}000$ the solve is $5.0$ to
$5.6$\,s of a $5.8$ to $6.5$\,s total, once compilation is charged once
rather than per rung, so the iteration count already dominates; at
$220{,}000$ the solve is $44$\,s of a $53$\,s total. On portfolio at
$50{,}223$, by contrast, compilation alone was $71.4\%$ of the total and the
solve was $8.1\%$. The decomposition $T = N \times t$ therefore has no single
dominant term across families or across the size range, and which factor to
attack depends on where the instance sits.

\paragraph{What governs the payoff, and what does not.}
It is tempting to read the rule as ``reduce $\|L\|$''. Two controlled
experiments say otherwise, each holding $\|L\|$ fixed by construction. In
the first, a ridge $\rho$ is swept on the $x$ block of a least-squares
problem with a dense knapsack row, which leaves $\|L\|$ identical at
$42.34 \to 3.87$ for every $\rho$; the median relocation gain over three
seeds moves from $3.19\times$ at $\rho=0$ to $102.5\times$ at $\rho=10^{-1}$.
In the second, the factor count $k$ of portfolio is swept at fixed $n$,
which changes the rank of the operator left behind; at $n=20{,}000$ the
relocated iteration count grows from $20$ at $k=10$ to $605$ at $k=2{,}000$,
about $k^{0.64}$, while $\|L\|$ after relocation moves only from $6.27$ to
$8.59$. So the payoff is set by the curvature the smooth term retains and by
the rank of the residual operator, and in neither sweep by $\|L\|$ itself.
In the same sweep the standard formulation fails to reach $10^{-6}$ in two
million iterations at every rank, so for this family relocation is not an
acceleration but the difference between solving and not solving.

\paragraph{Limits.}
The two quantities above predict our iteration count, not the speedup. The
speedup also depends on how hard the instance is for the baseline, and the
two families differ sharply there: Clarabel needs $9.06$\,s on the SVM dual
at $n_{\mathrm{var}} = 55{,}000$ and $0.253$\,s on portfolio at $50{,}223$.
The SVM dual wins at the smaller size because a direct solver struggles
there, not because the transform works better; we take $3{,}456$ iterations
on it against $256$ on portfolio. Two entries above are bounds rather than
measurements, since the standard arm did not converge within the cap. At
$n_{\mathrm{var}} = 550{,}000$ Clarabel did not return within one hour, so
those instances carry no certified reference and are excluded.

\paragraph{Negative control.}
A rule that fires everywhere is not a diagnosis, so the procedure has to be
able to say no, and for a structural reason rather than a threshold. At
$n_{\mathrm{var}}$ near $20{,}000$ per family:

\begin{center}
\begin{tabular}{lrrll}
\toprule
family & top row mass & singular gap & detector & plan when forced \\
\midrule
lasso     & 0.4454 & 1.0  & silent & no bounds on the relocated block \\
svm       & 0.0862 & 1.0  & silent & selected row is not an equality \\
huber     & 0.0647 & 1.0  & silent & only $21{,}671$ of $51{,}671$ bounded \\
control   & 0.0272 & 1.0  & silent & bound rows do not match the block \\
portfolio & 0.9997 & 21.8 & fires  & relocatable, $141$ rows left in $L$ \\
svm dual  & 0.9999 & 31.5 & fires  & relocatable, $2{,}000$ rows left in $L$ \\
\bottomrule
\end{tabular}
\end{center}

The separation is not delicate. The singular gap is exactly $1.0$ on all
four families we already win and between $21.8$ and $31.5$ on the two that
need the transform, and the top row mass separates $0.027$ to $0.445$ from
$0.9997$ to $0.9999$. The second column of refusals matters as much as the
first: when the detector is overridden and the transform is attempted
anyway, it still declines on all four, each time because a stated structural
condition fails rather than because a constant was tuned.

\section{What honest accounting costs, and the corrected table}
\label{app:accounting}

Two defects in our own measurement were found and fixed on September 14, and
every number in this paper that was measured through the adaptive driver is
affected by both. They pull in opposite directions, so the corrected figures
are not a uniform rescaling of the old ones.

\paragraph{Charging the tolerance ladder.}
Our solver walks a ladder of solver tolerances
$(\varepsilon, \varepsilon/10, \varepsilon/100, \varepsilon/1000)$ and
returns at the first rung whose answer passes acceptance on the original
problem. The harness charged wall-clock for the accepted rung only. That is
not what a user pays: the solver cannot know in advance which tolerance will
certify, so walking the ladder is the work, and every rung must be charged.
The correction is largest exactly where acceptance happens late. On the SVM
dual every instance accepts at the last rung, so four rungs run and one was
billed; on huber acceptance is at the second rung. Both harnesses now record
each rung with its iteration count and timing, which is also what makes the
failure analysis below possible.

\paragraph{A compile cache.}
Profiling the ladder showed compilation being paid on every rung: $1.011$\,s
then $0.938$\,s on huber, $51\%$ of that solve's wall clock. The cause was
not the tolerance, which never enters the compiled function and is tested in
Python on the returned metric. It was that the chunk function is a fresh
closure on every call of the driver, so the JIT sees a new function object
and misses its own cache although the computation is byte-identical. Caching
the compiled executable against the built scheme removes the redundant
compilations. Correctness was checked rather than assumed: iteration counts
and the accepted rung are identical, and the objective differs by
$1.3\times10^{-16}$, one unit in the last place, which is the same spread
seen between two repetitions within a single process and is therefore
floating-point non-determinism in the GPU reductions rather than an effect of
the cache. A first version of the cache held the built scheme alive without
bound, which pinned GPU memory for every instance it had seen and produced
out-of-memory failures; it is now bounded, and memory returns to baseline
after every family.

\paragraph{The corrected table.}
All figures below are on one build with both corrections, at tolerance
$10^{-6}$, three seeds, with intervals over seeds. Each arm pays its own
setup on every repetition and baselines run in separate processes.

\begin{center}
\begin{tabular}{llrrl}
\toprule
family & $n_{\mathrm{var}}$ & as first reported & corrected & interval \\
\midrule
huber     & 140{,}000 & $63.89\times$ & $45.66\times$ & $[40.75, 51.17]$ \\
lasso     & 100{,}000 & $55.60\times$ & $49.57\times$ & $[36.77, 66.83]$ \\
svm       & 110{,}000 & $4.17\times$  & $3.78\times$  & $[3.67, 3.90]$ \\
control   & 16{,}200  & $1.896\times$ & $1.27\times$  & $[1.07, 1.50]$ \\
portfolio & 250{,}500 & $1.89\times$  & $1.57\times$  & $[1.28, 1.92]$ \\
portfolio & 500{,}707 & $3.09\times$  & $2.37\times$  & $[1.78, 3.16]$ \\
svm dual  & 55{,}000  & $3.51\times$  & $1.45\times$  & $[1.26, 1.67]$ \\
svm dual  & 220{,}000 & $8.97\times$  & $5.29\times$  & $[5.09, 5.50]$ \\
svm dual  & 550{,}000 & $17.79\times$ & $7.91\times$  & $[6.34, 9.88]$ \\
MEB dual  & 44{,}000  & --            & $4.71\times$  & $[3.48, 6.37]$ \\
MEB dual  & 165{,}000 & --            & $46.81\times$ & $[26.19, 83.67]$ \\
\bottomrule
\end{tabular}
\end{center}

The table is complete: every family we report now has a figure on this one
build. Every family still wins and every interval excludes $1.0$. The
corrections cost between a sixth and three fifths of the original figure
depending on how late acceptance happens, and they changed no conclusion.
Control is the narrowest of them at $1.27\times$ $[1.07, 1.50]$, down from
$1.896\times$, and it is the family where the honest accounting bites hardest
because its objective gap sits at $9.5\times10^{-7}$ against a $10^{-6}$ bar
and it therefore needs the full ladder to certify. Control is measured on the
H100 because it cannot be run on the other machine: at
$n_{\mathrm{var}} = 16{,}200$ the compiler constant-folds an
eighteen-million-element gather and exhausts the memory left beside a
resident inference server. The two MEB dual rows have no first-reported
figure because that family was added after the corrections were in place and
has only ever been measured this way.

\paragraph{Crossovers move with the accounting.}
Both families lose to a direct solver below a crossover and win above it, and
correcting the accounting moved both crossovers upward. The SVM dual crosses
between $n_{\mathrm{var}} = 38{,}500$ ($0.79\times$, $[0.72, 0.87]$) and
$55{,}000$ ($1.45\times$); portfolio crosses between $100{,}316$
($0.92\times$, $[0.79, 1.07]$) and $250{,}500$ ($1.57\times$). Reporting a
single speedup for either family without its size would be meaningless.

\paragraph{A failure mode worth reporting.}
At $n_{\mathrm{var}} = 11{,}000$ two of three SVM dual seeds did not certify.
The per-rung log shows why, and it is not the iteration cap: the run used
$1{,}024$ to $4{,}224$ iterations of a $200{,}000$ budget, and the row
violation was still falling across the ladder
($7.1\times10^{-5}$, $1.6\times10^{-5}$, $5.7\times10^{-6}$,
$3.5\times10^{-6}$) when the ladder ran out at $10^{-9}$. A longer ladder
would very likely certify, cheaply. The same log exposes a related fact: the
driver's internal certificate under-reports true feasibility on the relocated
form by one to two orders, declaring convergence at target $10^{-6}$ while
the true row violation on the original problem is $7.1\times10^{-5}$. Nothing
incorrect was accepted, because acceptance is judged on the original problem
and not on that certificate, but it explains why acceptance lands late and
why the ladder is needed at all.

\paragraph{Why this is reported at all.}
A solver paper is a measurement paper, and every defect above inflated our
own numbers or masked a failure. We report them because the alternative,
quietly publishing the larger figures, is the failure mode that makes this
literature hard to build on. The corrected numbers are smaller and we have
more confidence in them.

\section{The claim this appendix supports}
\label{app:claim}

Stated plainly, so it can be attacked: in a primal-dual first-order method,
which constraints live in the linear operator and which live in the proximal
step is a first-class design choice; it can be decided automatically from the
problem data; and it reaches a point that the preconditioners used in
practice do not reach and that the optimal diagonal preconditioner reaches
only by solving an auxiliary optimization problem. The three clauses are
separable and each is supported by a different measurement. The third clause
was originally written as ``it is not reducible to preconditioning'' and that
stronger form is false; Section~\ref{app:priorart} reports the experiment
that refuted it and the weaker statement that survived.

\paragraph{It is a design choice, not a detail.}
On portfolio at $n_{\mathrm{var}} = 20{,}141$ the standard formulation does
not reach $10^{-6}$ in two million iterations at any factor rank we tried,
while the relocated formulation needs between $20$ and $605$. That is not an
acceleration, it is the difference between solving and not solving, and the
two formulations describe the same feasible set. The choice of where a
constraint lives is therefore not a numerical detail but a decision that
determines whether the method works at all on that family.

\paragraph{It can be decided from the data.}
The diagnosis is a two-line test on the rescaled operator: does one row carry
most of the top left singular direction, and is the operator effectively rank
one at the top. Across six families the separation is not delicate. The
singular gap is exactly $1.0$ on the four families that need no transform and
between $21.8$ and $31.5$ on the two that do; the top row mass separates
$0.027$ to $0.445$ from $0.9997$ to $0.9999$. The transform itself is then
read off $(A, l, u)$ rather than written per family, and it declines when any
structural precondition fails, which is what it does on all four of the
others even when the detector is overridden. No constant in the procedure was
tuned to make this separation appear.

\paragraph{It is not the preconditioning anyone runs.}
This is the clause most likely to be challenged, since a dense row is what a
diagonal preconditioner exists to handle, and the obvious objection is that
one could simply scale the row. One can: scaling row $j$ by $\varepsilon$
gives an equivalent problem with $\|L\|$ as small as desired, and a fully
optimized pair of diagonals does make the iteration converge, as
Section~\ref{app:priorart} reports. What is true is narrower. The best
hand-tuned single-row scaling on portfolio still needs $101{,}500$ iterations
where relocation needs at most $50$; the response is not monotone in
$\varepsilon$, so the constant would have to be searched rather than derived;
Ruiz equilibration, the preconditioner actually in use, is invariant to the
entire sweep; and the optimal diagonal preconditioner is itself the solution of
an auxiliary optimization and lands at about $1{,}500$ to $1{,}750$ iterations.
This paragraph originally said that relocation needs $250$ iterations and that
Ruiz converges at no value of $\varepsilon$. Both were wrong: $250$ was the
checkpoint stride of the loop, not a measured count, and Ruiz does converge on
portfolio once given the wall clock the optimized diagonal spends (at $334{,}510$
to $434{,}010$ iterations across three seeds), which the original $200{,}000$
cap hid. The seeded, wall-time comparison is in Appendix~\ref{app:field}: the
optimized diagonal loses to Ruiz in wall time, and relocation beats both on this
family. The mechanism is that the dual step is a single global
scalar while the dense row and the sparse rows want different ones;
relocation removes that trade-off by taking the row out of the dual problem,
where a diagonal scaling can only rebalance it. The hand-tuned and Ruiz parts
of this shape reproduce on the SVM dual; the optimized-diagonal comparison has
so far been run on portfolio only.

\paragraph{A third family, and the other projection branch.}
The rule was supported by two families, one of them added in the same week,
which is thin. The third is the dual of the minimum enclosing ball, the
Chebyshev centre of a point set, whose feasible set is exactly the
probability simplex. It is classical, its provenance is computational
geometry rather than finance or machine learning, and as with the SVM dual
the Gram matrix is dense and is never formed: the centre is lifted as a
variable so the coupling stays sparse. We checked that it is the minimum
enclosing ball and not merely a quadratic program of the right shape, by
recovering the centre from the optimal multipliers and confirming that no
random perturbation of it reduces the maximum distance to the points.

The diagnosis fires on it with no family-specific code, at row mass $0.9998$
to $0.9999$ and singular gap $24.2$ to $36.2$, and the derived plan selects
the simplex branch rather than the box-hyperplane branch, which is the first
time that branch was chosen for a family it was not written for. Its speedup over
Clarabel, three seeds at each size under the accounting of
Section~\ref{app:accounting}, is

\begin{center}
\begin{tabular}{lrl}
\toprule
$n_{\mathrm{var}}$ & speedup & $95\%$ interval \\
\midrule
5{,}500   & $0.02\times$  & $[0.01, 0.03]$ \\
44{,}000  & $4.71\times$  & $[3.48, 6.37]$ \\
165{,}000 & $46.81\times$ & $[26.19, 83.67]$ \\
\bottomrule
\end{tabular}
\end{center}

which is the steepest growth of the three relocation families: at the largest
size the relocated form needs $1{,}216$ iterations and $3.3$ to $3.6$\,s
against a direct solver's $140$ to $210$\,s. The interval is wide because the
per-seed figures range from $38.2\times$ to $60.4\times$, which is itself
worth noting: the margin at this size is large enough that seed-to-seed
variation in the baseline's factorization dominates the spread.

The three families now span both regimes the mechanisms distinguish.
Portfolio's smooth term is strictly curved on the decision block; the SVM
dual and the minimum enclosing ball dual both have a zero block there, so
$\mu/L = 0$ exactly. Two use the simplex projection and one the
box-hyperplane projection. The rule was not modified for any of them.

\paragraph{What governs the payoff.}
Two quantities of the post-relocation problem, and $\|L\|$ is neither. A
ridge sweep that holds $\|L\|$ fixed by construction moves the gain from
$3.19\times$ to $102.5\times$ as the curvature retained by the smooth term
grows; a rank sweep that holds curvature fixed moves the relocated iteration
count from $20$ to $605$ while $\|L\|$ moves by a factor of $1.37$. Both
predict our iteration count rather than the speedup, which additionally
depends on the instance's difficulty for the baseline and on where the
instance sits relative to the setup-dominance crossover.

\paragraph{What this claim is not.}
It is not a claim that a search over a parameter family found a better
setting; the decision is discrete and structural, and the same search space
contains no scalar that reproduces it. It is not a claim of a new algorithm:
the iteration is Condat-Vu throughout, and an operator search over scheme
variants found nothing beating it at realistic size. It is not a claim that
the transform is always available, since it requires an exact projection onto
the relocated set, which exists for the simplex and for a box intersected
with a hyperplane but not in general. And the third clause is supported by
measurement on two families, not by a theorem; we state the mechanism, we do
not prove that no diagonal scaling can succeed.

\section{The largest open risk: prior art on the rule itself}
\label{app:priorart}

The single largest threat to this paper is not a measurement. It is that the
diagnose-and-relocate rule, or something close enough to it, is already
published. That question is unresolved as of September 14, our own search
having been inconclusive, and it should be settled before the main body is
written, because it determines what the paper is allowed to claim rather than
merely how the claim is phrased.

\paragraph{What is certainly not new.}
The transform itself is not a contribution and we do not present it as one.
Exact projection onto the probability simplex is classical and we use the
standard method; projection onto a box intersected with a hyperplane is the
continuous knapsack problem and is equally standard. Putting a simplex
projection inside the proximal step of a primal-dual method is routine
practice in entropic optimal transport, where the constraint structure makes
it the obvious formulation, and nobody would call that new. The literature on
where to split a problem between smooth, proximal and linear terms is large
and old, and the observation that different splittings of the same problem
converge at different rates is not ours.

\paragraph{What might be new, in decreasing order of confidence.}
First, the negative result that the equilibration used in practice does not
substitute for relocation: that the best hand-tuned row scaling is still two
orders of magnitude worse than moving the row, that the response is not
monotone in the scaling, that Ruiz equilibration is invariant to the entire
sweep and converges at no value of it, and that the optimal diagonal
preconditioner closes the gap only partly and only by solving an auxiliary
problem. We have not seen this stated. Second, the two
quantities that govern the payoff, the curvature the smooth term retains and
the rank of the residual operator, each isolated by a sweep that holds the
other and the operator norm fixed, together with the negative result that the
operator norm predicts neither. Third, the diagnosis itself, a spectral test
on the rescaled operator that selects which families need the transform and
declines on the rest, and which is what makes the choice automatic rather
than a modelling decision made by hand.

\paragraph{What the contribution collapses to under each scenario.}
If the diagnosis is published, what remains is the mechanism and the
preconditioning refutation, which is a workshop-scale result and not an ICLR
paper on its own. If the mechanism is also published, nothing remains but a
solver engineering report with an honest measurement protocol. If only the
transform is published, which is the most likely case given the optimal
transport literature, then the diagnosis, the mechanism and the refutation
all stand and the paper is unaffected, because we already concede the
transform. The cost of being wrong about this is highest at the point where
it is cheapest to check.

\paragraph{What a first targeted search found, September 14.}
Searching for the diagnosis and the refutation rather than the transform
turned up the nearest prior art immediately, and it is classical. Van der
Sluis (1969) bounds how far simple scalings are from optimal: for a matrix
with at most $q$ nonzeros in any row, row normalization is within a factor of
$q$ of the best achievable diagonal scaling. That bound is vacuous precisely
when a row is dense, which is the case we care about, so the classical theory
already says dense rows are where diagonal preconditioning's guarantee is
weakest. It does not say that no diagonal scaling works there. There is
also an active line computing the optimal diagonal preconditioner explicitly,
posing it as a quasiconvex or semidefinite problem and solving it at scale
(Qu, Gao and Ye and successors, with scalable approximations reported for
matrices up to $200{,}000$), alongside older equilibration work.

This is a direct and fair challenge to the third clause as we currently
support it. Our refutation swept Ruiz equilibration and hand-tuned row
scaling; it did not sweep the optimal diagonal preconditioner. A reader is
entitled to ask whether the optimum, rather than the heuristic, closes the
gap. Until that is answered the third clause should be read as a statement
about the preconditioners actually used in practice by first-order solvers,
which is a weaker and still useful claim, and not as a statement about all
diagonal scalings. The experiment that settles it is small and is described
in the next paragraph.

\paragraph{What the targeted searches returned on the diagnosis.}
The searches that bear on the load-bearing claim, which is now the diagnosis
rather than the refutation, were for automatic or adaptive selection of the
splitting itself. They return a large and mature literature on selecting
something else: step sizes, penalty parameters, and above all the metric,
including diagonal scaling for Douglas-Rachford and ADMM and the metric
selection results that optimize linear convergence bounds. We did not find a
published criterion, spectral or otherwise, for deciding which constraints
should live in the proximal step rather than in the linear operator. That is
a negative result from a handful of searches and not a proof of absence, and
it should be repeated against the operator-splitting literature directly
before the main body relies on it.

One result deserves specific mention because it is the same problem family.
Moehle, Gindi, Boyd and Kochenderfer formulate portfolio construction as a
separable-affine problem and solve it by ADMM, using the same lifting we do,
introducing the factor exposure as a variable and writing the budget
constraint alongside it. Their split is the opposite of ours: the budget
constraint stays in the affine operator and only the separable terms are
handled by proximal operators, and their section on scaling lists diagonal
row and column scaling as the equivalence available to the formulation. Their
structure could not express the transform we apply, because projection onto
the simplex couples coordinates and their proximal terms are separable by
construction. We read this as support for the claim that the choice is not
obvious rather than as evidence against novelty: a leading group working on
exactly this problem chose the other split and treated scaling, not
splitting, as the free variable.

\paragraph{The nearest work in this venue, and why we are not it.}
The closest prior work at a machine learning venue is RLQP (Ichnowski, Jain,
Stellato, Banjac, Luo, Borrelli, Gonzalez, Stoica and Goldberg,
``Accelerating Quadratic Optimization with Reinforcement Learning'',
Advances in Neural Information Processing Systems 34, NeurIPS 2021,
pp.~21043--21055; arXiv:2107.10847). It trains a policy with the TD3
actor-critic method to adapt the penalty parameter $\rho$ of the OSQP
solver while it runs, either as a single scalar mapped to every constraint
or, through a multi-agent formulation sharing one policy, as a separate
coefficient for each constraint. It reports speedups of up to $3\times$ over
OSQP on OSQP's own benchmark classes, $1.30\times$ on the Netlib linear
programs, and a shifted geometric mean $1.829\times$ in its favour on the
Maros and M\'esz\'aros problems both solvers finish.

It is the right paper to define ourselves against for three reasons. It was
accepted, so solver acceleration is in scope for these venues. It evaluates on
the same benchmark suite: the control, Huber, SVM, lasso and portfolio
generators behind every family in Section~\ref{app:log:qp2026} are the
OSQP benchmark classes RLQP trains and tests on. And it is the strongest
existing instance of the category a reader will reach for, learning or tuning
a solver's parameters.

The difference is not one of degree. RLQP selects a continuous parameter
inside a fixed formulation; which constraints are dualized never changes, and
cannot, because the policy acts only on $\rho$. What we select is the
formulation itself, and the choice is discrete: a constraint lives either in
the linear operator or in the proximal step, with no interpolation between
the two. Section~\ref{app:priorart} reports that a fully optimized pair of
diagonal scalings, a far larger continuous family than one penalty vector,
still reaches a worse point on all three relocation families, so tuning every
continuous parameter we have would not find this. The two are orthogonal
rather than competing, and in principle compose: one could relocate first and
then adapt the penalty of whatever splitting remains.

Several of RLQP's own findings bear directly on ours. Its related work notes
that computing the optimal $\rho$ in closed form requires semidefinite
programs harder than the quadratic program itself, which is the same
observation we make about the optimal diagonal preconditioner, an auxiliary
optimization more expensive than the instance it is meant to accelerate. Its
policy degrades on problem dimensions outside its training range, required
several days of training on a 256\,GB machine, and in some cases lost its
iteration savings to the cost of evaluating the policy; the diagnosis here
needs no training and is recomputed from each instance's data. And its
experiments ran on CPU at dimensions up to roughly two thousand variables,
with portfolio instances between $5$ and $154$, where ours run on GPU at
$16{,}200$ to $550{,}000$.

One point of scope follows from how OSQP works and should be stated rather
than left for a reviewer. OSQP takes each step by solving a linear system in
the KKT matrix exactly, by factorization, so a dense constraint row enters it
through that factorization and not through a step-size condition on the norm
of the constraint operator. The pathology we diagnose is a property of
matrix-free primal-dual iterations whose admissible steps are set by $\|L\|$,
and it is consistent with RLQP obtaining gains on portfolio by adapting $\rho$
alone. We therefore claim the diagnosis for matrix-free first-order methods,
which are the methods that run on accelerators at the sizes above, and not for
factorization-based splitting solvers.

\paragraph{How it should be settled.}
Not by a keyword search for the transform, which will return the optimal
transport work and say nothing. The searches that matter are for the
diagnosis and the refutation: spectral criteria for choosing an operator
splitting; any statement that a dense constraint row resists diagonal
equilibration; adaptive or automatic splitting selection in primal-dual or
operator-splitting methods; and the preconditioning literature for
first-order conic and quadratic solvers, where a negative result of this shape
would most naturally have appeared. Until that is done, the claims in
Section~\ref{app:claim} should be read as claims about what we measured, not
as claims of priority.

\paragraph{We ran that experiment, and the third clause is false as stated.}
On portfolio at $n_{\mathrm{var}} = 828$ we optimized the full pair of
diagonals by L-BFGS with analytic gradients through the singular values,
minimizing the condition number of $D_1 A D_2$, which is the objective the
optimal-preconditioning literature targets. A diagonal scaling is an exact
reformulation here, since every bound in this family is a row of $A$.

\begin{center}
\begin{tabular}{lrrrr}
\toprule
scaling & $\kappa$ & $\sigma_1/\sigma_2$ & $\|L\|$ & iterations \\
\midrule
none                       & $76.27$   & $10.50$ & $28.87$ & $>200{,}000$ \\
Ruiz, 10 sweeps            & $48.76$   & $9.26$  & $15.86$ & $>200{,}000$ \\
optimized for $\kappa$      & $1.216$   & $1.00$  & $1.23$  & $1{,}750$ \\
optimized for $\sigma_1/\sigma_2$ & $457.0$ & $1.00$ & $1.78$ & $>200{,}000$ \\
relocated                  & --        & --      & $2.58$  & $250$ \\
\bottomrule
\end{tabular}
\end{center}

The optimized diagonal preconditioner converges. It drives the condition
number to $1.216$ and the singular gap to $1.00$, and the iteration that
never finished under Ruiz finishes in $1{,}750$ steps. So relocation is not
categorically beyond the reach of diagonal scaling, and any claim that it is
should be struck. What survives is narrower and we state only this: Ruiz
equilibration, the preconditioner first-order solvers actually use, does not
find such a scaling at any point of a hundredfold sweep; the optimum that
does is the solution of an auxiliary optimization problem, which here took
$82$\,s of L-BFGS on dense factorizations for a problem a direct solver
finishes in well under a second, and which at realistic sizes is the
quasiconvex or semidefinite problem of the cited literature; and relocation
reaches a better point, $250$ iterations against $1{,}750$, at no cost beyond
the diagnosis.

\paragraph{A diagnosis is not an objective.}
The fourth row of that table is instructive and we did not anticipate it.
Optimizing the detector's own statistic, the ratio $\sigma_1/\sigma_2$,
drives it to exactly $1.00$, flattening the spectral signature the detector
keys on, while the condition number degrades to $457$ and the iteration never
converges. The quantity that identifies which family needs a transform is not
the quantity to minimize. The detector locates the pathology; it does not
measure the remedy.

\paragraph{Which objective to precondition for inverts between families.}
Repeating the whole experiment on the SVM dual at $n_{\mathrm{var}} = 880$
gives a picture that is not a weaker version of the portfolio one but a
different one:

\begin{center}
\begin{tabular}{lrrrr}
\toprule
scaling & $\kappa$ & $\sigma_1/\sigma_2$ & $\|L\|$ & iterations \\
\midrule
none                       & $71.89$ & $11.14$ & $28.87$ & $146{,}750$ \\
Ruiz, 10 sweeps            & $68.34$ & $12.97$ & $28.81$ & $142{,}000$ \\
optimized for $\kappa$      & $1.466$ & $1.00$  & $1.31$  & $57{,}500$ \\
optimized for $\sigma_1/\sigma_2$ & $12.39$ & $1.00$ & $5.93$ & $4{,}500$ \\
relocated                  & --      & --      & $2.42$  & $1{,}000$ \\
\bottomrule
\end{tabular}
\end{center}

Running the same comparison on all three relocation families gives:

\begin{center}
\begin{tabular}{lrrrr}
\toprule
family & opt.\ for $\kappa$ & opt.\ for $\sigma_1/\sigma_2$ & relocated & ratio \\
\midrule
portfolio & $1{,}750$  & $>200{,}000$ & $250$   & $7\times$ \\
SVM dual  & $57{,}500$ & $4{,}500$    & $1{,}000$ & $4.5\times$ \\
MEB dual  & $39{,}500$ & $>200{,}000$ & $250$   & $158\times$ \\
\bottomrule
\end{tabular}
\end{center}

Minimizing the condition number rescues portfolio and the minimum enclosing
ball dual, both from non-convergence, and is nearly useless on the SVM dual,
which improves only from $142{,}000$ to $57{,}500$. Minimizing the singular
gap rescues the SVM dual and is worse than doing nothing on the other two.
The inversion is therefore specific to one family rather than a general
alternation, but it is real, and the consequence is the same: no single
scalar objective serves all three, and choosing between them would require
knowing in advance which regime an instance is in. Relocation requires no
such choice and wins on all three, by factors of $7$, $4.5$ and $158$ over
the best scaling we could find for each.

This is the honest form of the third clause and it is stronger than the
categorical version we withdrew, because it is a statement about what
preconditioning would have to know rather than about what it can achieve. It
also sharpens the practical point: the optimal scaling is the solution of an
auxiliary optimization, and the objective of that optimization is itself
instance-dependent.

One further observation belongs here because it is uncomfortable for the
standard pipeline rather than for our claim. On the minimum enclosing ball
dual, ten sweeps of Ruiz make the problem measurably worse, taking the
condition number from $73.26$ to $89.23$ and the singular gap from $10.93$ to
$14.40$. The equilibration that every first-order solver applies by default
is not merely unhelpful on this family; it moves in the wrong direction.

A caveat on scope. The standard formulation of the SVM dual does converge,
in $146{,}750$ iterations, where portfolio's does not converge at all. The
claim that relocation is the difference between solving and not solving is
therefore specific to portfolio; on the SVM dual it is a speedup of two
orders of magnitude over the unrelocated form and of four and a half over
the best diagonal scaling we could find.

\section{A proposed abstract for the constraint-placement thesis}
\label{app:abstract}

The main text of this submission argues an earlier thesis. This subsection
records what the abstract would say if the paper were rebuilt around the
result in Section~\ref{app:claim}, so that the framing can be judged against
the evidence without waiting on the rewrite. It is a proposal and nothing in
the main body depends on it.

\paragraph{Draft.}
In primal-dual first-order methods for convex quadratic programs, every
constraint must be placed either in the linear operator, where it is handled
by a dual variable, or in the proximal step, where it is handled by an exact
projection. This placement is usually fixed when the problem is modelled and
is not treated as a decision. We show that it is a decision, that it can be
made automatically from the problem data, and that it dominates the
parameters practitioners actually tune. A spectral test on the equilibrated
constraint operator identifies the cases where placement is wrong: across six
problem families the test separates cleanly, with a top singular gap of
exactly $1.0$ on the four families that need no change and between $21.8$ and
$36.2$ on the two that do, and the transform is then derived from the problem
data rather than written per family. On three families where it fires, moving
the identified constraint into the proximal step turns a formulation that
does not converge in two million iterations into one that converges in tens,
and yields end-to-end speedups over a production interior-point solver of
$2.4\times$ to $46.8\times$ at accuracy matched on the original problem. The
payoff is governed by the curvature the smooth term retains and by the rank
of the residual operator, and by neither the operator norm nor any quantity a
diagonal preconditioner optimizes: fully optimized diagonal scalings, a far
larger continuous family than any penalty parameter, reach a worse point on
all three families, and no single scalar objective serves all of them. We
report the measurement protocol that makes these numbers comparable, together
with the corrections it forced on our own earlier figures.

\paragraph{What the draft deliberately does not claim.}
It does not claim a new algorithm; the iteration is Condat-Vu throughout. It
does not claim the transform is new, only the diagnosis, the mechanism and
the comparison. It does not say ``not reducible to preconditioning'', which
we tested and retracted. And it does not quote a single speedup per family
without its size, because every family loses below a crossover.

\section{Planned experiments: four problem classes and the regimes where they should win}
\label{app:log:classes}

Written 2026-09-24, before any of the first three were run, so that the
predictions sit on the record ahead of the measurements. Each class states the
reason we expect to win, the regime where the win should appear, the datasets,
and the field that has to be completed for the comparison to count. Confidence
is a prediction, not a result, except for Class 4, which is already measured.

\subsection{Class 1. Simplex and box constrained least squares at scale, via
hyperspectral unmixing}

\paragraph{Confidence: highest of the three unrun classes.} It is the family we
already win, moved onto real data.

\paragraph{The problem.} Fully constrained least squares per pixel,
$\min \tfrac12\|Ax-b\|^2$ subject to $x \ge 0$ and $\mathbf{1}^\top x = 1$,
solved for every pixel in a scene.

\paragraph{Why we should win.} It is \texttt{knapsack\_ridge} with a real
application attached, and that family is our strongest surviving result
($4.1\times$ to $5.0\times$ at $n = 12{,}000$, $20.6\times$ to $35.0\times$ at
$48{,}000$, and no engine answering at all at $100{,}000$). The constraint is a
box intersected with a hyperplane, which is exactly the operator already
certified in this work at $37.35\,\mu$s per call against $117.88\,\mu$s for the
textbook breakpoint search. The relocation gene moves the simplex out of $L$
into $g$ and the synthesized projection is what makes it cheap, so this would
give the operator-synthesis contribution its first demonstration on real rather
than generated instances, which the paper's conclusion currently lists as
future work. The specialists are real but structurally doomed at scale:
\texttt{scipy.optimize.nnls} is Lawson-Hanson active set and single threaded,
and SUnSAL is a per-pixel ADMM loop. That gives the shape of curve the
specialist clause needs, deferring correctly on one pixel and winning by orders
of magnitude on the scene.

\paragraph{Hardware role.} $A$ is small and dense, bands by endmembers, but the
batch is about $50{,}000$ pixels, so the iteration is a batched GEMM plus a
batched projection. The device gene therefore fires on batch size rather than on
instance size, which is a second and independent confirmation of the
$\mathrm{nnz}(L) + n$ ruler. Precision should fire as well: float32 is expected
to certify $10^{-4}$ and to stall before $10^{-8}$.

\paragraph{Datasets.} AVIRIS Cuprite, $250 \times 191$ at 188 usable bands,
about $47{,}750$ simultaneous simplex-constrained problems, which is the
standard unmixing benchmark. Urban HYDICE, $307 \times 307$ at 162 bands and a
different endmember count, as the transfer check.

\paragraph{Field to complete.} \texttt{scipy.optimize.nnls},
\texttt{scipy.optimize.lsq\_linear}, SUnSAL and FCLS from \texttt{pysptools},
CVXPY to Clarabel and SCS on a pixel subset, and cuML where it applies.

\subsection{Class 2. Isotropic total variation inverse problems, 2D at high
resolution and 3D volumes}

\paragraph{Confidence: high on the hardware claim, medium on beating a well
tuned primal-dual baseline.}

\paragraph{The problem.}
$\min \tfrac12\|Ax-y\|^2 + \lambda\,\mathrm{TV}_{\mathrm{iso}}(x)$ with $x$ in a
pixel box, for $A = I$ (denoising) and $A$ a blur (deblurring).

\paragraph{Why we should win.} This is the one class where the fusion tile is
native. The temporal kernel runs $T$ iterations per launch on a $(B+2T)^2$ tile
with halo in shared memory, which is a stencil, and total variation's $L$ is the
discrete gradient. Measured on that kernel: $1.15\times$ to $1.66\times$ over
the XLA loop on the RTX 4090 and $1.95\times$ to $2.74\times$ on the H100, with
the best tile differing at five of six sizes. Hardware selection is the
mechanism here rather than a decoration. It also repairs the objection that
killed the earlier TV cell: anisotropic TV-L2 is exactly solvable by parametric
max-flow, which is why that comparison field was empty, but parametric max-flow
is anisotropic only. Isotropic TV-L2 has no exact combinatorial solver and TV
deblurring has none at all, so the field can be completed honestly, with
max-flow included on an anisotropic control to show the router deferring. The
pixel box is a free exact projection that most implementations simply drop, and
including it exactly is a formulation win available to us and not to them.
Finally, the specialists are themselves Chambolle-Pock and FISTA, which live
inside our genome space, so we compete on evolved stepsizes, certified
over-relaxation, the fused kernel and the device, all of which they hand-set.

\paragraph{Predicted regime.} For an $m \times m$ image,
$\mathrm{nnz}(L) + n \approx 5m^2$, so the CPU-to-GPU crossover at about
$22{,}400$ lands near $m = 67$. A sweep from $32^2$ to $2048^2$ therefore
crosses the ruler inside the sweep, which is a clean out-of-family test of the
crossover claim.

\paragraph{Datasets.} DIV2K, 800 images at about $2040 \times 1356$ and no
registration required, used as the resolution sweep from $64^2$ to full 2K. A 3D
CT or OCT volume stack, $512 \times 512 \times 200$ and $n \approx 5\times10^7$,
as the extreme, where CVXPY-class solvers cannot form the problem at all and
only a fused GPU first-order method survives.

\paragraph{Field to complete.}
\texttt{skimage.restoration.denoise\_tv\_chambolle}, a tuned Chambolle-Pock,
FISTA on the dual, ProxImaL, CVXPY to SCS and Clarabel at small sizes, and
parametric max-flow on the anisotropic control.

\subsection{Class 3. Wasserstein distributionally robust learning}

\paragraph{Confidence: high against the field, and the least de-risked of the
three.}

\paragraph{The problem.} Distributionally robust logistic or linear regression
over a Wasserstein ball, whose convex reformulation carries one constraint block
per training sample.

\paragraph{Why we should win.} There is no specialist at all. Practitioners
write the problem in CVXPY and call Clarabel, SCS or Mosek, so under this
paper's field-completion protocol this is the cleanest available instance of
outperforming general-purpose baselines where no specialist exists, and the one
class where completing the field costs us nothing. The per-sample blocks are
norm balls and simplices with exact projections, so the relocation gene applies
directly, and the problem is block separable, so the iteration is embarrassingly
parallel. Problem size grows linearly in the sample count, so interior-point
meets the wall it met on \texttt{knapsack\_ridge} at $n = 100{,}000$, where no
engine returned an answer. It is also the most legible of the four to an ICLR
audience.

\paragraph{Datasets.} Covertype, $581{,}012$ samples and 54 features, large
enough that CVXPY cannot canonicalize the full problem, with a well-defined
small subsample where Clarabel wins and the router should defer. Adult,
$48{,}842$ samples, as the small control that establishes the crossover rather
than only the win.

\paragraph{Caveat on hardware.} This class exercises device admission, the
$\mathrm{nnz}(L) + n$ ruler and the thread-width gene, but not the fusion tile,
because there is no stencil. If all classes must exercise the tile, the swap is
compressed-sensing MRI on fastMRI, which combines an FFT operator, a total
variation stencil and a decisive float32 regime. It is not ranked higher only
because BART is a strong and battle-tested baseline.

\subsection{Class 4. The lasso regularization path, by exact homotopy}

\paragraph{Confidence: the highest of the four, because this one is already
measured rather than predicted.} It is also the reference derivation of the
rediscovery argument in the paper's methods: the four decisive steps below are
what the loop is asked to re-derive as four single mutations.

\paragraph{The problem.} Solutions of
$P_\lambda(w) = \tfrac{1}{2n}\|y - Xw\|_2^2 + \lambda\|w\|_1$ at a decreasing
grid $\lambda_1 > \dots > \lambda_K$.

\paragraph{Why we win.} The path is piecewise linear in $\lambda$, so an exact
traversal obtains every requested grid point inside a segment by interpolation
rather than by solving. Against iterative competitors the margin therefore grows
with the accuracy target instead of shrinking, from $1.72\times$ at $10^{-4}$ to
$5.15\times$ at $10^{-12}$, which is the signature of an exact method against
iterative ones. Reachability is the claim to defend first: the engine certifies
$349$ of $350$ (dataset, target) cells where the best of fifteen baselines
certifies $290$, and every baseline has a region of the accuracy range it cannot
enter at any tolerance setting. Four implementation decisions carry the result.
The factorization of $X_A^\top X_A$ is maintained across the whole path by
Cholesky update and downdate rather than refactorized, drifting
$9.5\times10^{-16}$ over 300 interleaved insert and delete operations, so the
cheap version is also the accurate one. The screening bound is free because it
reuses $G_{AA}d = s$, the equation that defines the direction, and it is exact,
asserted by comparing step counts rather than objectives. The traversal hands
its tail to a working-set engine at a predicate derived from writing out both
costs, where $n$ and $p$ cancel and the decision reduces to the active-set size
against a multiple of the grid size. And that active-set size is obtained by
counting variables as they enter rather than by bounding the support with
$\min(n,p)$: replacing the bound with the measurement was worth $2.40\times$ to
$3.33\times$, larger than any kernel change in the work.

\paragraph{Hardware role.} The decision to use a device is taken before touching
the CUDA stack, because importing the runtime and creating a context costs more
than a small solve; the quantity asked for is the bytes the traversal will
stream, not the size of the design. The width gene appears at CPU granularity
rather than as a tile: a homotopy step is a handful of small BLAS calls, so a
path issues thousands of tiny parallel regions and pays a barrier for each, and
a $500 \times 1000$ path measured $2.55$\,ms on 16 threads against
$261.84$\,ms on 64, a $103\times$ regression. There is no fusion tile here,
because a homotopy step is not a stencil.

\paragraph{Datasets.} The whole LIBSVM catalog that loads as a single-target
regression and fits the host's free disk: 70 datasets spanning
\texttt{iris} at $150 \times 4$ to \texttt{finance} at
$16{,}087 \times 4{,}265{,}669$, at five accuracy targets each. Separately,
SimpleTES's own 17-problem benchmark under SimpleTES's own protocol and validity
rule, which supplies a comparison of one discovery system against another on the
same problem.

\paragraph{Field completed.} Fifteen solvers over six algorithm classes:
scikit-learn's LARS for exact homotopy; celer, skglm, scikit-learn, glmnet,
SimpleTES, glum and statsmodels for coordinate descent; FISTA for proximal
gradient; Clarabel and HiGHS for interior point; SCS and OSQP for operator
splitting; and cuML on the device. Three further libraries are named as
unavailable rather than omitted.

\paragraph{What has to happen before this counts as our result.} The numbers
were measured on the OptiVault host at a median load average of $15.8$ on a
shared 64-core machine, with a worst single-cell drift of $4.50\times$ between
passes, so they are claimed as ratios within a dataset and never as absolute
times. To enter the paper they have to be re-run on the RTX 4090 and the H100,
and the engine has to be registered as a routable engine behind the existing
permutation-invariant detector, so that the $428.6\times$ sparse-lasso loss
reported in the paper's limits section is re-measured rather than argued away.

\subsection{How the four divide the work}

\begin{center}\small
\begin{tabular}{@{}p{2.05cm}p{2.65cm}p{2.65cm}p{1.85cm}p{2.65cm}p{2.15cm}@{}}
\toprule
class & specialist situation & relocation & fusion tile & device ruler &
precision gene \\
\midrule
1 constrained LS & exists, dies at scale & the certified operator & no &
yes, through batch size & yes \\
\addlinespace
2 isotropic TV & none exact & the pixel box & yes, native & yes, crossover
inside the sweep & yes \\
\addlinespace
3 Wasserstein DRO & none at all & per-sample blocks & no & yes & partly \\
\addlinespace
4 lasso path & fifteen, three win somewhere & not needed, the family choice
does it & no & yes, on streamed bytes & no, f64 throughout \\
\bottomrule
\end{tabular}
\end{center}

Class 1 gives operator synthesis its first demonstration on real data. Class 2
is the only one where the fusion tile is the actual mechanism. Class 3 is the
cleanest field-completion win and the most legible to the venue. Class 4 is the
only one already measured, and it is the class that supplies the evidence the
specialist clause of the abstract currently lacks. If one is to be run first it
should be Cuprite, because it reuses an operator already certified here, its
baselines install without friction, and it should produce a result within a day
on the RTX 4090.


\subsection{Baseline policy: naive, out of the box}
\label{app:log:baselinepolicy}

Every competitor is run in its published or packaged form, as shipped. We do
not rewrite a competitor's inner loop, swap its linear algebra, or hand
optimize it. The reason is not politeness. A reviewer can check an
out-of-the-box baseline against its own repository; nobody can check our
private improvement of someone else's code, and ``we made their solver faster
and still beat it'' only invites an argument about whether we made it fast
enough.

Two things are explicitly NOT covered by that rule, because both favor the
competitor. Sweeping a solver's own exposed knobs, its tolerance and its
iteration cap, is required by the protocol of
Section~\ref{PAPER-sec:exp:protocol} and is how a competitor is given its
best legitimate setting. And using a library the way its documentation says
to, for instance amortizing CVXPY's canonicalization across instances with a
parameter rather than rebuilding the problem per pixel, is ordinary correct
usage rather than tuning.

\paragraph{Audit of what was done before the policy was fixed.} Every
departure is recorded here with its direction, because the dangerous direction
is the one that makes a baseline slower than default and none of these do.

\begin{center}\small
\begin{tabular}{@{}p{3.3cm}p{5.2cm}p{2.3cm}p{4.0cm}@{}}
\toprule
baseline & what was changed & effect on the baseline & disposition \\
\midrule
SUnSAL, unmixing &
published \texttt{cho\_solve} replaced by an explicit inverse, faster
projection, residuals every 25 iterations &
faster, $1.74\times$ ($13.3 \to 23.1$ it/s) &
\textbf{reverted}; headline re-run in the published form, the faster form kept
as a disclosed control \\
\addlinespace
CVXPY Clarabel and ECOS, unmixing &
a \texttt{cp.Parameter} so canonicalization amortizes across pixels &
faster & kept, this is documented CVXPY usage \\
\addlinespace
scikit-image TV, denoising &
used as shipped; its own \texttt{eps} and \texttt{max\_num\_iter} swept &
faster, we take its fastest passing setting & kept, protocol \\
\addlinespace
CVXPY SCS, deblurring &
sparse blur operator and \texttt{cp.tv} rather than a dense matrix &
faster; dense would exhaust memory at $128^2$ & kept, this is how the problem
would be written \\
\bottomrule
\end{tabular}
\end{center}

Because every one of those made the competitor faster, every speedup reported
for the three new classes is conservative: the out-of-the-box numbers can only
move in our favor. The one case where a choice of ours had made a competitor
look worse than it is was not a baseline change at all but a grading error,
scoring scikit-image by a dual point constructed from its own iterate, and it
was corrected before any number was reported (Section~\ref{app:log:classes}).

\section{The hardware layer is a headline, not a postscript}
\label{app:log:hardware}

Recorded 2026-09-24, alongside the three new problem classes. The claim this
section supports is that hardware selection is a \textbf{gene}, evaluated inside
the search and coupled back into which algorithm wins, rather than a tuning
pass applied to an already chosen algorithm.

\subsection{Three genes, one quantity}

Layer 8 is not a post-processing step applied to the winning algorithm. It is
one of three hardware genes, all set from the same quantity, the work per step
of whatever inner loop layers 1 through 6 produced.

\begin{center}\small
\begin{tabular}{@{}p{1.5cm}p{2.3cm}p{3.5cm}p{2.3cm}p{4.3cm}@{}}
\toprule
gene & question & GPU form & CPU form & the cliff that proves it is a gene \\
\midrule
7 admission & cross the device at all? &
streamed bytes against the bandwidth difference, decided before touching the
runtime & stay &
importing the runtime costs more than a small solve \\
\addlinespace
8 width & how much work per launch or barrier? &
fusion tile: $T$ iterations per launch on a $(B{+}2T)^2$ tile with halo in
shared memory & BLAS thread count from work per step &
the wrong device's tile costs $1.06\times$ to $1.63\times$; $64$ threads
instead of $16$ costs $103\times$ \\
\addlinespace
6 contract & what precision, and where &
per-expression fastmath scope & same &
fastmath on a Schur complement turned a $10^{-14}$ certificate into
$10^{-2}$ \\
\bottomrule
\end{tabular}
\end{center}

The fusion tile and the thread count are the same gene at two granularities.
Both answer how wide the parallel region should be relative to the work the
algorithm gives it, both are read off the device profile, and both have a
cliff when set by an environment default instead of by the algorithm. That is
what makes the RTX 4090 against H100 tile result and the $16$ against $64$
thread result one finding rather than two anecdotes.

\subsection{What the three new classes add}

\paragraph{Gene 6 fires hardest, and the cliff is a wrong answer.}
On the Cuprite library, whose condition number is $7.5\times10^8$, float32
does not merely stall. It diverges. The worst-pixel relative gap ends at
$5.26$ on the RTX 4090 and $7.595$ on the H100, and \textbf{nothing is
certified at any target on either machine}, while float64 certifies $10^{-4}$
in $165.3$\,s and $53.3$\,s respectively. A solver trusting its own stopping
rule would have returned those iterates. The acceptance check refuses them.
Precision is therefore not a speed knob on this class; it is an admissibility
condition, and it is device independent, which is itself worth recording
because the other two genes are not.

\paragraph{Gene 7, and a confirmation the cost model needed.}
The same arm on the same instance took \textbf{identical} iteration counts on
both devices, $8{,}501$ to reach $10^{-3}$ and $16{,}001$ to reach $10^{-4}$,
with final gaps agreeing to fourteen digits
($9.999381424056729\times10^{-6}$ on the H100 against
$9.999381424162939\times10^{-6}$ on the RTX 4090). Only the wall time differs,
by $3.10\times$ at $10^{-4}$. The cost model assumes $N$ is device independent
and prices only $t(\text{operator},\text{device})$; this is an out-of-family
confirmation of that assumption on a class the model has never seen.

\paragraph{Gene 8, reproduced, and currently unexploited.}
The temporally fused TV kernel was re-run on 2026-09-24 and reproduces the
September selector exactly, picking the same tile at every size:

\begin{center}\small
\begin{tabular}{@{}lrrrrr@{}}
\toprule
 & \multicolumn{2}{c}{RTX 4090 (17 tiles admissible)} &
   \multicolumn{2}{c}{H100 (24 admissible)} & \\
\cmidrule(lr){2-3}\cmidrule(lr){4-5}
size & tile $(B,T)$ & gain & tile $(B,T)$ & gain & devices agree? \\
\midrule
$128$ & $(8,4)$  & $1.619\times$ & $(16,8)$ & $2.521\times$ & no \\
$256$ & $(24,4)$ & $1.482\times$ & $(16,8)$ & $2.330\times$ & no \\
$384$ & $(24,4)$ & $1.270\times$ & $(24,8)$ & $1.883\times$ & no \\
$496$ & $(32,4)$ & $1.337\times$ & $(32,8)$ & $1.792\times$ & no \\
\bottomrule
\end{tabular}
\end{center}

The devices disagree at every size, and they disagree in a structured way:
\textbf{the RTX 4090 never selects $T = 8$ and the H100 selects it at every
size.} The reason is visible in the device profile rather than in the
timings. The 4090 offers $101{,}376$ bytes of opt-in shared memory per block
and admits only 17 of the 21 tiles; the H100 offers $232{,}448$ and admits all
of them, so it can afford to buy HBM traffic with $2.25\times$ to $4.00\times$
halo recomputation that the 4090's rate-limited consumer float64 cannot. This
is the same pattern the September sweep recorded, reproduced on both machines
on 2026-09-24 with the same tile selected at every one of the eight cells.

Against the XLA compiled loop the two effects compound: at $128$ the fused
kernel is about $3.25\times$ faster than XLA, because the hand-written single
step is already $2.01\times$ faster and fusion adds $1.619\times$ on top.

\textbf{This is the honest gap in the new numbers.} The denoising and
deblurring arms reported in Section~\ref{app:log:classes} run the XLA loop,
not the fused kernel, so every time reported there is an unfused time and the
fused kernel is a multiplier that has not been taken. The deblurring arm
cannot take it at all as written, because its $f$ is applied by FFT rather
than by a stencil, so only the TV term is fusable. Wiring the selector into
the denoising arm is the obvious next step and it is not done.

\subsection{The operator choice, and a measurement that had to be corrected}

The simplex projection was measured against a safeguarded bisection on the
dual scalar, nine repetitions, median reported with the spread:

\begin{center}\small
\begin{tabular}{@{}lrrr@{}}
\toprule
device & sort & bisection (48 steps) & winner \\
\midrule
RTX 4090 & $0.289$\,ms (spread $1.02\times$) & $0.795$\,ms & sort, $2.75\times$ \\
H100     & $0.107$\,ms (spread $1.01\times$) & $0.393$\,ms & sort, $3.67\times$ \\
CPU      & $13.3$\,ms & $58.3$\,ms & sort, $4.38\times$ \\
\bottomrule
\end{tabular}
\end{center}

An earlier single-shot reading of exactly this quantity reported the opposite
ordering on the H100 and was written up as a device-dependent flip of the
operator gene, with a $6.56\times$ cross-device penalty. It was wrong. One of
the two timings was taken while a full-scene job still held the RTX 4090, and
neither was repeated. Under repetition the ordering is the same on all three
devices and there is no flip. The benchmark now records load average and
spread beside every number, and the bisection variant is recorded in the
solver as measured and declined rather than left as an assumption. The
relevant card is \texttt{record\_the\_machine\_state\_beside\_the\_timing},
and it cost a retraction to relearn.

\subsection{Why this is hardware-aware and not hardware-tuned}

The measured $t(\text{operator},\text{device})$ feeds back into the layer-5
predicate. A better fused tile lowers $t$ for the iterative tail, which moves
the crossover, which changes which algorithm family wins on a given shape. The
hardware layer is not downstream of the algorithm layer; the two are coupled
through the cost model of Equation~\ref{PAPER-eq:costmodel}, and the loop
closes through that coupling. That single arc is the thesis: hardware-aware
loop discovery of optimization algorithms, where the device changes the
algorithm and not only its constants.

On the unmixing class that coupling is visible rather than hypothetical. The
first arm built there was FISTA, which needed more than $10^5$ iterations
because it cannot use the one factorization every pixel shares. Relocating the
simplex into $g$ changed the family to ADMM on a shared inverse and cut that
to thousands. Which family wins depends on the per-iteration price of the
shared solve, and that price is a device measurement.


\bibliographystyle{iclr2027/iclr2027_conference}
\bibliography{refs}

\end{document}
